%
% Dette dokument bruges til at "samle" rapporten ud fra de individuelle sektioner i .\sections\
%

\documentclass[11pt,twoside,a4paper,openright]{report}

%%
%% Preamble er en slags side med indstilligner for hele rapporten
%%
%% !!! Hvis du compiler rapporten imens denne fil er åben får du en fejl. Gå til master.tex for at compile !!!
%%


\documentclass[11pt,twoside,a4paper,openright]{report}
%%%%%%%%%%%%%%%%%%%%%%%%%%%%%%%%%%%%%%%%%%%%%%%%
% Language, Encoding and Fonts
% http://en.wikibooks.org/wiki/LaTeX/Internationalization
%%%%%%%%%%%%%%%%%%%%%%%%%%%%%%%%%%%%%%%%%%%%%%%%
% Select encoding of your inputs. Depends on
% your operating system and its default input
% encoding. Typically, you should use
%   Linux  : utf8 (most modern Linux distributions)
%            latin1 
%   Windows: ansinew
%            latin1 (works in most cases)
%   Mac    : applemac
% Notice that you can manually change the input
% encoding of your files by selecting "save as"
% an select the desired input encoding. 
\usepackage[utf8]{inputenc}
\usepackage{ifthen}
% Make latex understand and use the typographic
% rules of the language used in the document.
\usepackage[danish,english]{babel}
% Use the palatino font
\usepackage[sc]{mathpazo}
\linespread{1.05}         % Palatino needs more leading (space between lines)
% Choose the font encoding
\usepackage[T1]{fontenc}

%%%%%%%%%%%%%%%%%%%%%%%%%%%%%%%%%%%%%%%%%%%%%%%%
% Graphics and Tables
% http://en.wikibooks.org/wiki/LaTeX/Importing_Graphics
% http://en.wikibooks.org/wiki/LaTeX/Tables
% http://en.wikibooks.org/wiki/LaTeX/Colors
%%%%%%%%%%%%%%%%%%%%%%%%%%%%%%%%%%%%%%%%%%%%%%%%
% load a colour package
\usepackage{xcolor}
\definecolor{aaublue}{RGB}{33,26,82}% dark blue
% The standard graphics inclusion package
\usepackage{graphicx}
% Set up how figure and table captions are displayed
\usepackage{caption}
\captionsetup{%
  font=footnotesize,% set font size to footnotesize
  labelfont=bf % bold label (e.g., Figure 3.2) font
}
% Make the standard latex tables look so much better
\usepackage{array,booktabs}
% Enable the use of frames around, e.g., theorems
% The framed package is used in the example environment
\usepackage{framed}

% Adds support for full page background picture
\usepackage[contents={},color=gray]{background}
%\usepackage[contents=draft,color=gray]{background}

%%%%%%%%%%%%%%%%%%%%%%%%%%%%%%%%%%%%%%%%%%%%%%%%
% Mathematics
% http://en.wikibooks.org/wiki/LaTeX/Mathematics
%%%%%%%%%%%%%%%%%%%%%%%%%%%%%%%%%%%%%%%%%%%%%%%%
% Defines new environments such as equation,
% align and split 
\usepackage{amsmath}
% Adds new math symbols
\usepackage{amssymb}
% Use theorems in your document
% The ntheorem package is also used for the example environment
% When using thmmarks, amsmath must be an option as well. Otherwise \eqref doesn't work anymore.
\usepackage[framed,amsmath,thmmarks]{ntheorem}

%%%%%%%%%%%%%%%%%%%%%%%%%%%%%%%%%%%%%%%%%%%%%%%%
% Page Layout
% http://en.wikibooks.org/wiki/LaTeX/Page_Layout
%%%%%%%%%%%%%%%%%%%%%%%%%%%%%%%%%%%%%%%%%%%%%%%%
% Change margins, papersize, etc of the document
\usepackage[
  inner=28mm,% left margin on an odd page
  outer=41mm,% right margin on an odd page
  ]{geometry}
% Modify how \chapter, \section, etc. look
% The titlesec package is very configureable
\usepackage{titlesec}
\titleformat{\chapter}[display]{\normalfont\huge\bfseries}{\chaptertitlename\ \thechapter}{20pt}{\Huge}
\titleformat*{\section}{\normalfont\Large\bfseries}
\titleformat*{\subsection}{\normalfont\large\bfseries}
\titleformat*{\subsubsection}{\normalfont\normalsize\bfseries}
%\titleformat*{\paragraph}{\normalfont\normalsize\bfseries}
%\titleformat*{\subparagraph}{\normalfont\normalsize\bfseries}

% Clear empty pages between chapters
\let\origdoublepage\cleardoublepage
\newcommand{\clearemptydoublepage}{%
  \clearpage
  {\pagestyle{empty}\origdoublepage}%
}
\let\cleardoublepage\clearemptydoublepage

% Change the headers and footers
\usepackage{fancyhdr}
\pagestyle{fancy}
\fancyhf{} %delete everything
\renewcommand{\headrulewidth}{0pt} %remove the horizontal line in the header
\fancyhead[RE]{\small\nouppercase\leftmark} %even page - chapter title
\fancyhead[LO]{\small\nouppercase\rightmark} %uneven page - section title
\fancyhead[LE,RO]{\thepage} %page number on all pages
% Do not stretch the content of a page. Instead,
% insert white space at the bottom of the page
\raggedbottom
% Enable arithmetics with length. Useful when
% typesetting the layout.
\usepackage{calc}

%%%%%%%%%%%%%%%%%%%%%%%%%%%%%%%%%%%%%%%%%%%%%%%%
% Bibliography
% http://en.wikibooks.org/wiki/LaTeX/Bibliography_Management
%%%%%%%%%%%%%%%%%%%%%%%%%%%%%%%%%%%%%%%%%%%%%%%%
\usepackage[backend=biber,bibencoding=utf8]{biblatex}
\addbibresource{sections/6- bibliography.bib}
\usepackage{csquotes}
%%%%%%%%%%%%%%%%%%%%%%%%%%%%%%%%%%%%%%%%%%%%%%%%
% Misc
%%%%%%%%%%%%%%%%%%%%%%%%%%%%%%%%%%%%%%%%%%%%%%%%
% Add bibliography and index to the table of
% contents
\usepackage[nottoc]{tocbibind}
% Add the command \pageref{LastPage} which refers to the
% page number of the last page
\usepackage{lastpage}
% Add todo notes in the margin of the document
\usepackage[
%  disable, %turn off todonotes
  colorinlistoftodos, %enable a coloured square in the list of todos
  textwidth=\marginparwidth, %set the width of the todonotes
  textsize=scriptsize, %size of the text in the todonotes
  ]{todonotes}

%%%%%%%%%%%%%%%%%%%%%%%%%%%%%%%%%%%%%%%%%%%%%%%%
% Hyperlinks
% http://en.wikibooks.org/wiki/LaTeX/Hyperlinks
%%%%%%%%%%%%%%%%%%%%%%%%%%%%%%%%%%%%%%%%%%%%%%%%
% Enable hyperlinks and insert info into the pdf
% file. Hypperref should be loaded as one of the 
% last packages
\usepackage{hyperref}
\hypersetup{%
	%pdfpagelabels=true,%
	plainpages=false,%
	pdfauthor={Author(s)},%
	pdftitle={Title},%
	pdfsubject={Subject},%
	bookmarksnumbered=true,%
	colorlinks=false,%
	citecolor=black,%
	filecolor=black,%
	linkcolor=black,% you should probably change this to black before printing
	urlcolor=black,%
	pdfstartview=FitH%
}

%%
%% !!! Hvis du compiler rapporten imens denne fil er åben får du en fejl. Gå til master.tex for at compile !!!
%% % package inclusion and set up of the document
% see, e.g., http://en.wikibooks.org/wiki/LaTeX/Formatting#Hyphenation
% for more information on word hyphenation
\hyphenation{ex-am-ple hy-phen-a-tion short}
\hyphenation{long la-tex}
 % 
%%
%% Denne fil er til at lave macroer som kan kaldes i andre filer, hvis denne fil er defineret
%% Filen defineres med: %%
%% Denne fil er til at lave macroer som kan kaldes i andre filer, hvis denne fil er defineret
%% Filen defineres med: %%
%% Denne fil er til at lave macroer som kan kaldes i andre filer, hvis denne fil er defineret
%% Filen defineres med: \input{setup/macros.tex}
%%

%%%%%%%%%%%%%%%%%%%%%%%%%%%%%%%%%%%%%%%%%%%%%%%%
% Referencer
%%%%%%%%%%%%%%%%%%%%%%%%%%%%%%%%%%%%%%%%%%%%%%%%

% myPageref - Reference med angivelse af figur/tabel og sidetal. Figur/Tabel er med stort begyndelsesbogstav
\newcommand{\myAutoRefOnPage}[1]{\textit{\autoref{#1} on page \pageref{#1}}}
% Use: "this can be seen in \autorefonpage{fig:cool_graph}"
% Output: this can be seen in Figure 3.2 on page 23

%%%%%%%%%%%%%%%%%%%%%%%%%%%%%%%%%%%%%%%%%%%%%%%%
% Dato og tid
%%%%%%%%%%%%%%%%%%%%%%%%%%%%%%%%%%%%%%%%%%%%%%%%

% \ordinalth{<number>} will output st, nd, rd, or th
\newcommand{\ordinalTh}[1]{%
  \ifcase#1
    \textsuperscript{th}% case 0
  \or
    \textsuperscript{st}% case 1
  \or
    \textsuperscript{nd}% case 2
  \or
    \textsuperscript{rd}% case 3
  \else
    \textsuperscript{th}% everything else (4+)
  \fi
}

%%%%%%%%%%%%%%%%%%%%%%%%%%%%%%%%%%%%%%%%%%%%%%%%
% Misc-macroer
%%%%%%%%%%%%%%%%%%%%%%%%%%%%%%%%%%%%%%%%%%%%%%%%

% myTodayKort - Print af dags dato på kort format (DD/MMM-YYYY)
\newcommand{\myTodayKort}{\ifnum\number\day<10 0\fi \number\day /\ifcase \month \or jan\or feb\or mar\or apr\or may \or jun\or jul\or aug\or sep\or oct\or nov\or dec\fi -\number \year} 

% myTodayLang - Print af dags dato på langt format (DD. MMMM YYYY)(engelske måneder)
\newcommand{\myTodayLang}{\ifnum\number\day<10 0\fi \number\day \ordinalTh{\number\day} \space\ifcase \month \or January\or February\or March\or April\or May \or June\or July\or August\or September\or October\or November\or December\fi \space\number \year} 
% myTodayLangDA - Print af dags dato på langt format (DD. MMMM YYYY)(danske måneder)
\newcommand{\myTodayLangDA}{\ifnum\number\day<10 0\fi \number\day .\space\ifcase \month \or januar\or februar\or marts\or april\or maj \or juni\or juli\or august\or september\or oktober\or november\or december\fi \space\number \year} 

% myAppendixPageNumber - Nulstiller sidetal, så det specificeres for det enkelte bilag og formaterer sidetallet med prefix som sættes med parameteren.
\newcommand{\myAppendixPageNumbering}[1]{%
    \pagenumbering{arabic} % resets page-counter to 1
    \renewcommand*{\thepage}{{#1}\arabic{page}} 
}

% myTitle - Formatering, der gør teksten stor og fed
\newcommand{\myTitle}[1]{\noindent\textbf{\large{#1}}\par\noindent}

% myExpression - Formaterer teksten i kursiv og med dobbelte citationspile
\newcommand{\myExpression}[1]{»\textit{{#1}}«}

% mytodo - Macro til Todo der udfylder den firkantede parantes så man hurtigere kan skrive den - color bliver cyan, hvis andet parameter ikke er fyldt, textcolor bliver white, hvis tredje parameter ikke er udfyldt. Altså \mytodo{tekst}[color][textcolor]
% Eks1: \mytodo{det her er en cyan todo}
% Eks2: \mytodo{hej med dig}[red]
% Eks3: \mytodo{hej med dig med sort skrift}[green][black]
\usepackage{xparse}
\NewDocumentCommand{\mytodo}{m o o}{%
  \IfNoValueTF{#2}{\def\mytodocolor{cyan}}{\def\mytodocolor{#2}}%
  \IfNoValueTF{#3}{\def\mytodotextcolor{white}}{\def\mytodotextcolor{#3}}%
  \todo[inline, color=\mytodocolor, textcolor=\mytodotextcolor]{#1}%
}

%%%%%%%%%%%%%%%%%%%%%%%%%%%%%%%%%%%%%%%%%%%%%%%%
% Macros for the titlepage
%%%%%%%%%%%%%%%%%%%%%%%%%%%%%%%%%%%%%%%%%%%%%%%%
%Creates the aau titlepage
\newcommand{\aautitlepage}[3]{%
  {
    %set up various length
    \ifx\titlepageleftcolumnwidth\undefined
      \newlength{\titlepageleftcolumnwidth}
      \newlength{\titlepagerightcolumnwidth}
    \fi
    \setlength{\titlepageleftcolumnwidth}{0.48\textwidth-\tabcolsep}
    \setlength{\titlepagerightcolumnwidth}{\textwidth-2\tabcolsep-\titlepageleftcolumnwidth}
    %create title page
    \thispagestyle{empty}
    \noindent%
    \begin{tabular}{@{}ll@{}}
      \parbox{\titlepageleftcolumnwidth}{
        \iflanguage{danish}{%
          \includegraphics[width=\titlepageleftcolumnwidth-10mm]{media/AAUgraphics/aau_logo_da.pdf}
        }{%
          \includegraphics[width=\titlepageleftcolumnwidth-10mm]{media/AAUgraphics/aau_logo_en.pdf}
        }
      } &
      \parbox{\titlepagerightcolumnwidth}{\raggedleft\sf\small
        #2
      }\bigskip\\
       #1 &
      \parbox[t]{\titlepagerightcolumnwidth}{%
      \textbf{Abstract:}\bigskip\par
        \fbox{\parbox{\titlepagerightcolumnwidth-2\fboxsep-2\fboxrule}{%
          #3
        }}
      }\\
    \end{tabular}
    \vfill \bigskip
    \iflanguage{danish}{\noindent{\footnotesize\emph{Rapportens indhold er frit tilgængeligt, men offentliggørelse\\ (med kildeangivelse) må kun ske efter aftale med forfatterne.}}
    }{\noindent{\footnotesize\emph{The content of this report is freely available, but publication\\ (with reference) may only be pursued due to agreement with the author.}}
    }
    \clearpage
  }
}

%Create english project info
\newcommand{\englishprojectinfo}[8]{%
  \parbox[t]{\titlepageleftcolumnwidth}{
    \textbf{Title:}\\ #1\bigskip\par
    \textbf{Theme:}\\ #2\bigskip\par
    \textbf{Project Period:}\\ #3\bigskip\par
    \textbf{Project Group:}\\ #4\bigskip\par
    \textbf{Participant(s):}\\ #5\bigskip\par
    \textbf{Supervisor(s):}\\ #6\bigskip\par
    \textbf{Copies:} #7\bigskip\par
    \textbf{Number of pages:} \pageref{myLastPage}\bigskip\par
    \textbf{Date of Completion:}\\ #8
  }
}

%Create danish project info
\newcommand{\danishprojectinfo}[8]{%
  \parbox[t]{\titlepageleftcolumnwidth}{
    \textbf{Titel:}\\ #1\bigskip\par
    \textbf{Tema:}\\ #2\bigskip\par
    \textbf{Projektperiode:}\\ #3\bigskip\par
    \textbf{Projektgruppe:}\\ #4\bigskip\par
    \textbf{Deltager(e):}\\ #5\bigskip\par
    \textbf{Vejleder(e):}\\ #6\bigskip\par
    \textbf{Oplagstal:} #7\bigskip\par
    \textbf{Sidetal:} \pageref{myLastPage}\bigskip\par
    \textbf{Afleveringsdato:}\\ #8
  }
}

%%%%%%%%%%%%%%%%%%%%%%%%%%%%%%%%%%%%%%%%%%%%%%%%
% An example environment
%%%%%%%%%%%%%%%%%%%%%%%%%%%%%%%%%%%%%%%%%%%%%%%%
\theoremheaderfont{\normalfont\bfseries}
\theorembodyfont{\normalfont}
\theoremstyle{break}
\def\theoremframecommand{{\color{gray!50}\vrule width 5pt \hspace{5pt}}}
\newshadedtheorem{exa}{Example}[chapter]
\newenvironment{example}[1]{%
		\begin{exa}[#1]
}{%
		\end{exa}
}
%%

%%%%%%%%%%%%%%%%%%%%%%%%%%%%%%%%%%%%%%%%%%%%%%%%
% Referencer
%%%%%%%%%%%%%%%%%%%%%%%%%%%%%%%%%%%%%%%%%%%%%%%%

% myPageref - Reference med angivelse af figur/tabel og sidetal. Figur/Tabel er med stort begyndelsesbogstav
\newcommand{\myAutoRefOnPage}[1]{\textit{\autoref{#1} on page \pageref{#1}}}
% Use: "this can be seen in \autorefonpage{fig:cool_graph}"
% Output: this can be seen in Figure 3.2 on page 23

%%%%%%%%%%%%%%%%%%%%%%%%%%%%%%%%%%%%%%%%%%%%%%%%
% Dato og tid
%%%%%%%%%%%%%%%%%%%%%%%%%%%%%%%%%%%%%%%%%%%%%%%%

% \ordinalth{<number>} will output st, nd, rd, or th
\newcommand{\ordinalTh}[1]{%
  \ifcase#1
    \textsuperscript{th}% case 0
  \or
    \textsuperscript{st}% case 1
  \or
    \textsuperscript{nd}% case 2
  \or
    \textsuperscript{rd}% case 3
  \else
    \textsuperscript{th}% everything else (4+)
  \fi
}

%%%%%%%%%%%%%%%%%%%%%%%%%%%%%%%%%%%%%%%%%%%%%%%%
% Misc-macroer
%%%%%%%%%%%%%%%%%%%%%%%%%%%%%%%%%%%%%%%%%%%%%%%%

% myTodayKort - Print af dags dato på kort format (DD/MMM-YYYY)
\newcommand{\myTodayKort}{\ifnum\number\day<10 0\fi \number\day /\ifcase \month \or jan\or feb\or mar\or apr\or may \or jun\or jul\or aug\or sep\or oct\or nov\or dec\fi -\number \year} 

% myTodayLang - Print af dags dato på langt format (DD. MMMM YYYY)(engelske måneder)
\newcommand{\myTodayLang}{\ifnum\number\day<10 0\fi \number\day \ordinalTh{\number\day} \space\ifcase \month \or January\or February\or March\or April\or May \or June\or July\or August\or September\or October\or November\or December\fi \space\number \year} 
% myTodayLangDA - Print af dags dato på langt format (DD. MMMM YYYY)(danske måneder)
\newcommand{\myTodayLangDA}{\ifnum\number\day<10 0\fi \number\day .\space\ifcase \month \or januar\or februar\or marts\or april\or maj \or juni\or juli\or august\or september\or oktober\or november\or december\fi \space\number \year} 

% myAppendixPageNumber - Nulstiller sidetal, så det specificeres for det enkelte bilag og formaterer sidetallet med prefix som sættes med parameteren.
\newcommand{\myAppendixPageNumbering}[1]{%
    \pagenumbering{arabic} % resets page-counter to 1
    \renewcommand*{\thepage}{{#1}\arabic{page}} 
}

% myTitle - Formatering, der gør teksten stor og fed
\newcommand{\myTitle}[1]{\noindent\textbf{\large{#1}}\par\noindent}

% myExpression - Formaterer teksten i kursiv og med dobbelte citationspile
\newcommand{\myExpression}[1]{»\textit{{#1}}«}

% mytodo - Macro til Todo der udfylder den firkantede parantes så man hurtigere kan skrive den - color bliver cyan, hvis andet parameter ikke er fyldt, textcolor bliver white, hvis tredje parameter ikke er udfyldt. Altså \mytodo{tekst}[color][textcolor]
% Eks1: \mytodo{det her er en cyan todo}
% Eks2: \mytodo{hej med dig}[red]
% Eks3: \mytodo{hej med dig med sort skrift}[green][black]
\usepackage{xparse}
\NewDocumentCommand{\mytodo}{m o o}{%
  \IfNoValueTF{#2}{\def\mytodocolor{cyan}}{\def\mytodocolor{#2}}%
  \IfNoValueTF{#3}{\def\mytodotextcolor{white}}{\def\mytodotextcolor{#3}}%
  \todo[inline, color=\mytodocolor, textcolor=\mytodotextcolor]{#1}%
}

%%%%%%%%%%%%%%%%%%%%%%%%%%%%%%%%%%%%%%%%%%%%%%%%
% Macros for the titlepage
%%%%%%%%%%%%%%%%%%%%%%%%%%%%%%%%%%%%%%%%%%%%%%%%
%Creates the aau titlepage
\newcommand{\aautitlepage}[3]{%
  {
    %set up various length
    \ifx\titlepageleftcolumnwidth\undefined
      \newlength{\titlepageleftcolumnwidth}
      \newlength{\titlepagerightcolumnwidth}
    \fi
    \setlength{\titlepageleftcolumnwidth}{0.48\textwidth-\tabcolsep}
    \setlength{\titlepagerightcolumnwidth}{\textwidth-2\tabcolsep-\titlepageleftcolumnwidth}
    %create title page
    \thispagestyle{empty}
    \noindent%
    \begin{tabular}{@{}ll@{}}
      \parbox{\titlepageleftcolumnwidth}{
        \iflanguage{danish}{%
          \includegraphics[width=\titlepageleftcolumnwidth-10mm]{media/AAUgraphics/aau_logo_da.pdf}
        }{%
          \includegraphics[width=\titlepageleftcolumnwidth-10mm]{media/AAUgraphics/aau_logo_en.pdf}
        }
      } &
      \parbox{\titlepagerightcolumnwidth}{\raggedleft\sf\small
        #2
      }\bigskip\\
       #1 &
      \parbox[t]{\titlepagerightcolumnwidth}{%
      \textbf{Abstract:}\bigskip\par
        \fbox{\parbox{\titlepagerightcolumnwidth-2\fboxsep-2\fboxrule}{%
          #3
        }}
      }\\
    \end{tabular}
    \vfill \bigskip
    \iflanguage{danish}{\noindent{\footnotesize\emph{Rapportens indhold er frit tilgængeligt, men offentliggørelse\\ (med kildeangivelse) må kun ske efter aftale med forfatterne.}}
    }{\noindent{\footnotesize\emph{The content of this report is freely available, but publication\\ (with reference) may only be pursued due to agreement with the author.}}
    }
    \clearpage
  }
}

%Create english project info
\newcommand{\englishprojectinfo}[8]{%
  \parbox[t]{\titlepageleftcolumnwidth}{
    \textbf{Title:}\\ #1\bigskip\par
    \textbf{Theme:}\\ #2\bigskip\par
    \textbf{Project Period:}\\ #3\bigskip\par
    \textbf{Project Group:}\\ #4\bigskip\par
    \textbf{Participant(s):}\\ #5\bigskip\par
    \textbf{Supervisor(s):}\\ #6\bigskip\par
    \textbf{Copies:} #7\bigskip\par
    \textbf{Number of pages:} \pageref{myLastPage}\bigskip\par
    \textbf{Date of Completion:}\\ #8
  }
}

%Create danish project info
\newcommand{\danishprojectinfo}[8]{%
  \parbox[t]{\titlepageleftcolumnwidth}{
    \textbf{Titel:}\\ #1\bigskip\par
    \textbf{Tema:}\\ #2\bigskip\par
    \textbf{Projektperiode:}\\ #3\bigskip\par
    \textbf{Projektgruppe:}\\ #4\bigskip\par
    \textbf{Deltager(e):}\\ #5\bigskip\par
    \textbf{Vejleder(e):}\\ #6\bigskip\par
    \textbf{Oplagstal:} #7\bigskip\par
    \textbf{Sidetal:} \pageref{myLastPage}\bigskip\par
    \textbf{Afleveringsdato:}\\ #8
  }
}

%%%%%%%%%%%%%%%%%%%%%%%%%%%%%%%%%%%%%%%%%%%%%%%%
% An example environment
%%%%%%%%%%%%%%%%%%%%%%%%%%%%%%%%%%%%%%%%%%%%%%%%
\theoremheaderfont{\normalfont\bfseries}
\theorembodyfont{\normalfont}
\theoremstyle{break}
\def\theoremframecommand{{\color{gray!50}\vrule width 5pt \hspace{5pt}}}
\newshadedtheorem{exa}{Example}[chapter]
\newenvironment{example}[1]{%
		\begin{exa}[#1]
}{%
		\end{exa}
}
%%

%%%%%%%%%%%%%%%%%%%%%%%%%%%%%%%%%%%%%%%%%%%%%%%%
% Referencer
%%%%%%%%%%%%%%%%%%%%%%%%%%%%%%%%%%%%%%%%%%%%%%%%

% myPageref - Reference med angivelse af figur/tabel og sidetal. Figur/Tabel er med stort begyndelsesbogstav
\newcommand{\myAutoRefOnPage}[1]{\textit{\autoref{#1} on page \pageref{#1}}}
% Use: "this can be seen in \autorefonpage{fig:cool_graph}"
% Output: this can be seen in Figure 3.2 on page 23

%%%%%%%%%%%%%%%%%%%%%%%%%%%%%%%%%%%%%%%%%%%%%%%%
% Dato og tid
%%%%%%%%%%%%%%%%%%%%%%%%%%%%%%%%%%%%%%%%%%%%%%%%

% \ordinalth{<number>} will output st, nd, rd, or th
\newcommand{\ordinalTh}[1]{%
  \ifcase#1
    \textsuperscript{th}% case 0
  \or
    \textsuperscript{st}% case 1
  \or
    \textsuperscript{nd}% case 2
  \or
    \textsuperscript{rd}% case 3
  \else
    \textsuperscript{th}% everything else (4+)
  \fi
}

%%%%%%%%%%%%%%%%%%%%%%%%%%%%%%%%%%%%%%%%%%%%%%%%
% Misc-macroer
%%%%%%%%%%%%%%%%%%%%%%%%%%%%%%%%%%%%%%%%%%%%%%%%

% myTodayKort - Print af dags dato på kort format (DD/MMM-YYYY)
\newcommand{\myTodayKort}{\ifnum\number\day<10 0\fi \number\day /\ifcase \month \or jan\or feb\or mar\or apr\or may \or jun\or jul\or aug\or sep\or oct\or nov\or dec\fi -\number \year} 

% myTodayLang - Print af dags dato på langt format (DD. MMMM YYYY)(engelske måneder)
\newcommand{\myTodayLang}{\ifnum\number\day<10 0\fi \number\day \ordinalTh{\number\day} \space\ifcase \month \or January\or February\or March\or April\or May \or June\or July\or August\or September\or October\or November\or December\fi \space\number \year} 
% myTodayLangDA - Print af dags dato på langt format (DD. MMMM YYYY)(danske måneder)
\newcommand{\myTodayLangDA}{\ifnum\number\day<10 0\fi \number\day .\space\ifcase \month \or januar\or februar\or marts\or april\or maj \or juni\or juli\or august\or september\or oktober\or november\or december\fi \space\number \year} 

% myAppendixPageNumber - Nulstiller sidetal, så det specificeres for det enkelte bilag og formaterer sidetallet med prefix som sættes med parameteren.
\newcommand{\myAppendixPageNumbering}[1]{%
    \pagenumbering{arabic} % resets page-counter to 1
    \renewcommand*{\thepage}{{#1}\arabic{page}} 
}

% myTitle - Formatering, der gør teksten stor og fed
\newcommand{\myTitle}[1]{\noindent\textbf{\large{#1}}\par\noindent}

% myExpression - Formaterer teksten i kursiv og med dobbelte citationspile
\newcommand{\myExpression}[1]{»\textit{{#1}}«}

% mytodo - Macro til Todo der udfylder den firkantede parantes så man hurtigere kan skrive den - color bliver cyan, hvis andet parameter ikke er fyldt, textcolor bliver white, hvis tredje parameter ikke er udfyldt. Altså \mytodo{tekst}[color][textcolor]
% Eks1: \mytodo{det her er en cyan todo}
% Eks2: \mytodo{hej med dig}[red]
% Eks3: \mytodo{hej med dig med sort skrift}[green][black]
\usepackage{xparse}
\NewDocumentCommand{\mytodo}{m o o}{%
  \IfNoValueTF{#2}{\def\mytodocolor{cyan}}{\def\mytodocolor{#2}}%
  \IfNoValueTF{#3}{\def\mytodotextcolor{white}}{\def\mytodotextcolor{#3}}%
  \todo[inline, color=\mytodocolor, textcolor=\mytodotextcolor]{#1}%
}

%%%%%%%%%%%%%%%%%%%%%%%%%%%%%%%%%%%%%%%%%%%%%%%%
% Macros for the titlepage
%%%%%%%%%%%%%%%%%%%%%%%%%%%%%%%%%%%%%%%%%%%%%%%%
%Creates the aau titlepage
\newcommand{\aautitlepage}[3]{%
  {
    %set up various length
    \ifx\titlepageleftcolumnwidth\undefined
      \newlength{\titlepageleftcolumnwidth}
      \newlength{\titlepagerightcolumnwidth}
    \fi
    \setlength{\titlepageleftcolumnwidth}{0.48\textwidth-\tabcolsep}
    \setlength{\titlepagerightcolumnwidth}{\textwidth-2\tabcolsep-\titlepageleftcolumnwidth}
    %create title page
    \thispagestyle{empty}
    \noindent%
    \begin{tabular}{@{}ll@{}}
      \parbox{\titlepageleftcolumnwidth}{
        \iflanguage{danish}{%
          \includegraphics[width=\titlepageleftcolumnwidth-10mm]{media/AAUgraphics/aau_logo_da.pdf}
        }{%
          \includegraphics[width=\titlepageleftcolumnwidth-10mm]{media/AAUgraphics/aau_logo_en.pdf}
        }
      } &
      \parbox{\titlepagerightcolumnwidth}{\raggedleft\sf\small
        #2
      }\bigskip\\
       #1 &
      \parbox[t]{\titlepagerightcolumnwidth}{%
      \textbf{Abstract:}\bigskip\par
        \fbox{\parbox{\titlepagerightcolumnwidth-2\fboxsep-2\fboxrule}{%
          #3
        }}
      }\\
    \end{tabular}
    \vfill \bigskip
    \iflanguage{danish}{\noindent{\footnotesize\emph{Rapportens indhold er frit tilgængeligt, men offentliggørelse\\ (med kildeangivelse) må kun ske efter aftale med forfatterne.}}
    }{\noindent{\footnotesize\emph{The content of this report is freely available, but publication\\ (with reference) may only be pursued due to agreement with the author.}}
    }
    \clearpage
  }
}

%Create english project info
\newcommand{\englishprojectinfo}[8]{%
  \parbox[t]{\titlepageleftcolumnwidth}{
    \textbf{Title:}\\ #1\bigskip\par
    \textbf{Theme:}\\ #2\bigskip\par
    \textbf{Project Period:}\\ #3\bigskip\par
    \textbf{Project Group:}\\ #4\bigskip\par
    \textbf{Participant(s):}\\ #5\bigskip\par
    \textbf{Supervisor(s):}\\ #6\bigskip\par
    \textbf{Copies:} #7\bigskip\par
    \textbf{Number of pages:} \pageref{myLastPage}\bigskip\par
    \textbf{Date of Completion:}\\ #8
  }
}

%Create danish project info
\newcommand{\danishprojectinfo}[8]{%
  \parbox[t]{\titlepageleftcolumnwidth}{
    \textbf{Titel:}\\ #1\bigskip\par
    \textbf{Tema:}\\ #2\bigskip\par
    \textbf{Projektperiode:}\\ #3\bigskip\par
    \textbf{Projektgruppe:}\\ #4\bigskip\par
    \textbf{Deltager(e):}\\ #5\bigskip\par
    \textbf{Vejleder(e):}\\ #6\bigskip\par
    \textbf{Oplagstal:} #7\bigskip\par
    \textbf{Sidetal:} \pageref{myLastPage}\bigskip\par
    \textbf{Afleveringsdato:}\\ #8
  }
}

%%%%%%%%%%%%%%%%%%%%%%%%%%%%%%%%%%%%%%%%%%%%%%%%
% An example environment
%%%%%%%%%%%%%%%%%%%%%%%%%%%%%%%%%%%%%%%%%%%%%%%%
\theoremheaderfont{\normalfont\bfseries}
\theorembodyfont{\normalfont}
\theoremstyle{break}
\def\theoremframecommand{{\color{gray!50}\vrule width 5pt \hspace{5pt}}}
\newshadedtheorem{exa}{Example}[chapter]
\newenvironment{example}[1]{%
		\begin{exa}[#1]
}{%
		\end{exa}
} % Macroer der bruges igennem rapporten
%%
%% Denne fil bruges til at sætte variabler som bruges forskellige steder i rapporten
%% Hvis der er et felt man ikke ønsker at benytte skal det efterlades blankt
%% fx myDenneVariabelBrugesIkke{} - altså uden indhold mellem { }
%%
%% Det kan være nødvendigt at gennemgå de forskellige sider, for at fjerne referencer
%% til udkommenterede variabler, ellers kan resultatet se lidt underligt ud.
%%

%%% Formatering af sideopsætning
\def\myForsideMedBillede{1} % Skal forsiden være med billede? 1 = ja, 0 = nej
\def\myForsideBillede{media/AAUgraphics/frontpageimage.jpg} % Formateres i forholdet(højde:bredde) 1:2
\def\myPageOptimizeForPrint{1} % Skal dokumentet opsættes til print, hvor hvert nyt kapitel er på en lige side i forhold til starten? (Giver mange tomme sider)
% 1 = ja, opsæt til print / 0 = nej, opsæt til web/pdf
\def\myParagraphIndentDisable{1} % Skal indrykning efter afsnit deaktiveres?
% 1 = ja, 0 = nej
\def\myParagraphIndentLength{25pt} % Hvis nej, hvad skal indrykningen være?

% Sidemarginer kan desværre ikke opsættes her, de skal opsættes i "Preamble.tex"
% De findes under punktet 


%%% Tekst på forside/titlepages/colophon/preface
\def\myProjectTitle{Project Title: A Project}
\def\mySubtitle{A project that does sciency things}
\def\myProjectTheme{Theme of project}
\def\mySemester{Fall Semester} % 'Fall semester' eller 'Spring semester'
\def\mySemesterYear{2022}
\def\mySemesterNumber{1} % Hvilket semester er du på? (1,2,3,4,5,6,7,8,9,10)
\def\myProjectNumber{P0 Project} % Hvilket projekt er dette, fx P0 project eller Bachelor's Project
\def\myProjectECTS{15} % Hvor mange ECTS-point er projektet nomeret til?

% Tekst på danske title pages
\def\myProjectTitleDA{Projekttitel: Et projekt}
\def\myProjectThemeDA{Projekttema}
\def\mySemesterDA{Efterårssemester} % 'Forårssemesteret' eller 'Efterårssemesteret'

\def\myAuthorsNUM{4} % Hvor mange gruppemedlemmer der er (1-6)
% Herunder: Navne og mails på gruppemedlemmer
% hvis der fx kun er 4 i gruppen skal nummer 5 og 6 blot være tomme

\def\myAuthorOneName{Author 1} % Navn på gruppemedlem 1
\def\myAuthorTwoName{Author 2} % Navn på gruppemedlem 2
\def\myAuthorThreeName{Author 3} % Navn på gruppemedlem 3
\def\myAuthorFourName{Author 4} % Navn på gruppemedlem 4
\def\myAuthorFiveName{Author 5} % Navn på gruppemedlem 5
\def\myAuthorSixName{Author 6} % Navn på gruppemedlem 6

\def\myAuthorOneEmail{email-1@student.aau.dk} % Email på gruppemedlem 1
\def\myAuthorTwoEmail{email-2@student.aau.dk} % Email på gruppemedlem 2
\def\myAuthorThreeEmail{email-3@student.aau.dk} % Email på gruppemedlem 3
\def\myAuthorFourEmail{email-4@student.aau.dk} % Email på gruppemedlem 4
\def\myAuthorFiveEmail{email-5@student.aau.dk} % Email på gruppemedlem 5
\def\myAuthorSixEmail{email-6@student.aau.dk} % Email på gruppemedlem 6

\def\mySupervisorsNUM{2} % Hvor mange vejledere der er tilknyttet projektet? (1-3)
% Herunder: Navne på Vejledere
% hvis der fx kun er én vejleder tilknyttet projektet efterlades nummer 2 og 3 tomme

\def\mySupervisorOneName{Prof Supervisor 1} % Navn på vejleder 1
\def\mySupervisorOnePronoun{his} % Hvis der kun er én  vejleder, skriv his/her

\def\mySupervisorTwoName{Ass.Prof Supervisor 2} % Navn på vejleder 2
\def\mySupervisorThreeName{Supervisor 3} % Navn på vejleder 3

% Herunder lidt info om din uddannelse og dit institut
\def\myProgramme{Computer Engineering}
\def\myProgrammeDA{Computerteknologi}
\def\myGroupNumber{151}
\def\myDepartment{Department of Electronic Systems}
\def\myDepartmentDA{Institut for Elektroniske Systemer}
\def\myDepartmentWebsite{https://www.es.aau.dk}

% Colophon kommer til at stå på siden efter forsiden og er en reference
% til software og værktøjer brugt til at lave rapporten (og rapportens indhold)
\def\myColophon{%
Typesetting of this report has been done with Overleaf, a LaTeX writing tool.\\
\noindent
During the project simulations has been run in [instert software].\\
Figures have been created in [example: ]Lucidchart and Adobe Illustrator].
}

% Du er nu færdig med at udfylde DocumentSetup
% Kompiler rapporten og se om alt er som det skal være
% Fejl og mangler kan rapporteres som issues på GitHub
% https://github.com/Peasniped/AAU-LateX-Report-template/issues









%
%%%%% Nedenstående skal der ikke rettes i %%%%%
%

%% mySemesterNumber omdannes til Ordinaltal som bogstavtal og nummer (first og 1st)
\def\mySemesterNumberOrdinal{
\ifthenelse{\mySemesterNumber=1}{first}{
\ifthenelse{\mySemesterNumber=2}{second}{
\ifthenelse{\mySemesterNumber=3}{third}{
\ifthenelse{\mySemesterNumber=4}{fourth}{
\ifthenelse{\mySemesterNumber=5}{fifth}{
\ifthenelse{\mySemesterNumber=6}{sixth}{
\ifthenelse{\mySemesterNumber=7}{seventh}{
\ifthenelse{\mySemesterNumber=8}{eighth}{
\ifthenelse{\mySemesterNumber=9}{ninth}{
\ifthenelse{\mySemesterNumber=10}{tenth}{Fejl med def-mySemesterNumberOrdinal
}}}}}}}}}}}

\def\mySemesterNumberOrdinalNumberTh{
\ifthenelse{\mySemesterNumber=1}{1\textsuperscript{st}}{
\ifthenelse{\mySemesterNumber=2}{2\textsuperscript{nd}}{
\ifthenelse{\mySemesterNumber=3}{3\textsuperscript{rd}}{
\ifthenelse{\mySemesterNumber=4}{4\textsuperscript{th}}{
\ifthenelse{\mySemesterNumber=5}{5\textsuperscript{th}}{
\ifthenelse{\mySemesterNumber=6}{6\textsuperscript{th}}{
\ifthenelse{\mySemesterNumber=7}{7\textsuperscript{th}}{
\ifthenelse{\mySemesterNumber=8}{8\textsuperscript{th}}{
\ifthenelse{\mySemesterNumber=9}{9\textsuperscript{th}}{
\ifthenelse{\mySemesterNumber=10}{10\textsuperscript{th}}{Fejl med def-mySemesterNumberOrdinal
}}}}}}}}}}}

%% myAuthors dannes
\ifthenelse{\myAuthorsNUM=1}{
    \def\myAuthors{\myAuthorOneName}
}{ \ifthenelse{\myAuthorsNUM=2}{
    \def\myAuthors{\myAuthorOneName{}, \myAuthorTwoName}
}{ \ifthenelse{\myAuthorsNUM=3}{
    \def\myAuthors{\myAuthorOneName{}, \myAuthorTwoName{}, \myAuthorThreeName}
}{ \ifthenelse{\myAuthorsNUM=4}{
    \def\myAuthors{\myAuthorOneName{}, \myAuthorTwoName{}, \myAuthorThreeName{}, \myAuthorFourName}
}{ \ifthenelse{\myAuthorsNUM=5}{
    \def\myAuthors{\myAuthorOneName{}, \myAuthorTwoName{}, \myAuthorThreeName{}, \myAuthorFourName{}, \myAuthorFiveName}
}{ \ifthenelse{\myAuthorsNUM=6}{
    \def\myAuthors{\myAuthorOneName{}, \myAuthorTwoName{}, \myAuthorThreeName{}, \myAuthorFourName{}, \myAuthorFiveName{}, \myAuthorSixName}
}{ Fejl med def-myAuthors}}}}}}


%% myAuthorsNL dannes
\ifthenelse{\myAuthorsNUM=1}{
    \def\myAuthorsNL{\myAuthorOneName}
}{ \ifthenelse{\myAuthorsNUM=2}{
    \def\myAuthorsNL{\myAuthorOneName{},\\ \myAuthorTwoName}
}{ \ifthenelse{\myAuthorsNUM=3}{
    \def\myAuthorsNL{\myAuthorOneName{},\\ \myAuthorTwoName{},\\ \myAuthorThreeName}
}{ \ifthenelse{\myAuthorsNUM=4}{
    \def\myAuthorsNL{\myAuthorOneName{},\\ \myAuthorTwoName{},\\ \myAuthorThreeName{},\\ \myAuthorFourName}
}{ \ifthenelse{\myAuthorsNUM=5}{
    \def\myAuthorsNL{\myAuthorOneName{},\\ \myAuthorTwoName{},\\ \myAuthorThreeName{},\\ \myAuthorFourName{},\\ \myAuthorFiveName}
}{ \ifthenelse{\myAuthorsNUM=6}{
    \def\myAuthorsNL{\myAuthorOneName{},\\ \myAuthorTwoName{},\\ \myAuthorThreeName{},\\ \myAuthorFourName{},\\ \myAuthorFiveName{},\\ \myAuthorSixName}
}{ Fejl med def-myAuthorsNL}}}}}}


%% mySupervisorsNL dannes
\ifthenelse{\mySupervisorsNUM=1}{
    \def\mySupervisorsNL{\mySupervisorOneName}
}{ \ifthenelse{\mySupervisorsNUM=2}{
    \def\mySupervisorsNL{\mySupervisorOneName{},\\ \mySupervisorTwoName}
}{ \ifthenelse{\mySupervisorsNUM=3}{
    \def\mySupervisorsNL{\mySupervisorOneName{},\\ \mySupervisorTwoName{},\\ \mySupervisorThreeName}
}{ Fejl med def-mySupervisorsNL}}}


%% myThankOurSupervisors dannes
\ifthenelse{\mySupervisorsNUM=1}{
    \def\myThankOurSupervisors{extend our thanks to our supervisor \mySupervisorOneName{} for \mySupervisorOnePronoun{} help and guidance throughout the project.}    
}{ \ifthenelse{\mySupervisorsNUM=2}{
    \def\myThankOurSupervisors{extend our thanks to our supervisors \mySupervisorOneName{} and \mySupervisorTwoName{} for their help and guidance throughout the project.}
}{ \ifthenelse{\mySupervisorsNUM=3}{
    \def\myThankOurSupervisors{extend our thanks to our supervisors \mySupervisorOneName{}, \mySupervisorTwoName{} and \mySupervisorThreeName{} for their help and guidance throughout the project.}
}{ Fejl med def-myThankOurSupervisors}}}% Indlæs opsætning fra documentSetup
% ----- Her defineres forkortelser - de kaldes med \ac{example}

\DeclareAcronym{eksempel}{
    short =     {EKS},
    long =      {Eksempel},
    long-post = { (denne tekst skrives ved siden af forkortelsen første gang den bruges)},
    extra =     {Denne tekst vises i listen over forkortelser}
}


% ----- Her defineres begreber - de kaldes med \gls{example}

\newglossaryentry{eksempel}{
    name = {eksempel},
    description = {En eksemplificering af noget}
}


\makeglossaries % Indlæs definitioner på begreber og forkortelser

% Ret denne i documentSetup - Fjerner tomme sider mellem kapitler osv i hele dokumentet
\ifthenelse{\myPageOptimizeForPrint=1}{}{\let\cleardoublepage\clearpage}
% Ret denne i documentSetup - Fjerner Indrykning efter nyt afsnit
\ifthenelse{\myParagraphIndentDisable=1}{\setlength\parindent{0pt}}{\setlength\parindent{\myParagraphIndentLength}}

\begin{document}

\pagestyle{empty} %disable headers and footers
\pagenumbering{roman} %use roman page numbering in the frontmatter

% ToDo-oversigt - udkommenteres før at dokumentet afleveres
\listoftodos

\ifthenelse{\myForsideMedBillede=1}{%%
%% Denne fil bruges til at sætte variabler som bruges forskellige steder i rapporten
%% Hvis der er et felt man ikke ønsker at benytte skal det efterlades blankt
%% fx myDenneVariabelBrugesIkke{} - altså uden indhold mellem { }
%%
%% Det kan være nødvendigt at gennemgå de forskellige sider, for at fjerne referencer
%% til udkommenterede variabler, ellers kan resultatet se lidt underligt ud.
%%

%%% Formatering af sideopsætning
\def\myForsideMedBillede{1} % Skal forsiden være med billede? 1 = ja, 0 = nej
\def\myForsideBillede{media/AAUgraphics/frontpageimage.jpg} % Formateres i forholdet(højde:bredde) 1:2
\def\myPageOptimizeForPrint{1} % Skal dokumentet opsættes til print, hvor hvert nyt kapitel er på en lige side i forhold til starten? (Giver mange tomme sider)
% 1 = ja, opsæt til print / 0 = nej, opsæt til web/pdf
\def\myParagraphIndentDisable{1} % Skal indrykning efter afsnit deaktiveres?
% 1 = ja, 0 = nej
\def\myParagraphIndentLength{25pt} % Hvis nej, hvad skal indrykningen være?

% Sidemarginer kan desværre ikke opsættes her, de skal opsættes i "Preamble.tex"
% De findes under punktet 


%%% Tekst på forside/titlepages/colophon/preface
\def\myProjectTitle{Project Title: A Project}
\def\mySubtitle{A project that does sciency things}
\def\myProjectTheme{Theme of project}
\def\mySemester{Fall Semester} % 'Fall semester' eller 'Spring semester'
\def\mySemesterYear{2022}
\def\mySemesterNumber{1} % Hvilket semester er du på? (1,2,3,4,5,6,7,8,9,10)
\def\myProjectNumber{P0 Project} % Hvilket projekt er dette, fx P0 project eller Bachelor's Project
\def\myProjectECTS{15} % Hvor mange ECTS-point er projektet nomeret til?

% Tekst på danske title pages
\def\myProjectTitleDA{Projekttitel: Et projekt}
\def\myProjectThemeDA{Projekttema}
\def\mySemesterDA{Efterårssemester} % 'Forårssemesteret' eller 'Efterårssemesteret'

\def\myAuthorsNUM{4} % Hvor mange gruppemedlemmer der er (1-6)
% Herunder: Navne og mails på gruppemedlemmer
% hvis der fx kun er 4 i gruppen skal nummer 5 og 6 blot være tomme

\def\myAuthorOneName{Author 1} % Navn på gruppemedlem 1
\def\myAuthorTwoName{Author 2} % Navn på gruppemedlem 2
\def\myAuthorThreeName{Author 3} % Navn på gruppemedlem 3
\def\myAuthorFourName{Author 4} % Navn på gruppemedlem 4
\def\myAuthorFiveName{Author 5} % Navn på gruppemedlem 5
\def\myAuthorSixName{Author 6} % Navn på gruppemedlem 6

\def\myAuthorOneEmail{email-1@student.aau.dk} % Email på gruppemedlem 1
\def\myAuthorTwoEmail{email-2@student.aau.dk} % Email på gruppemedlem 2
\def\myAuthorThreeEmail{email-3@student.aau.dk} % Email på gruppemedlem 3
\def\myAuthorFourEmail{email-4@student.aau.dk} % Email på gruppemedlem 4
\def\myAuthorFiveEmail{email-5@student.aau.dk} % Email på gruppemedlem 5
\def\myAuthorSixEmail{email-6@student.aau.dk} % Email på gruppemedlem 6

\def\mySupervisorsNUM{2} % Hvor mange vejledere der er tilknyttet projektet? (1-3)
% Herunder: Navne på Vejledere
% hvis der fx kun er én vejleder tilknyttet projektet efterlades nummer 2 og 3 tomme

\def\mySupervisorOneName{Prof Supervisor 1} % Navn på vejleder 1
\def\mySupervisorOnePronoun{his} % Hvis der kun er én  vejleder, skriv his/her

\def\mySupervisorTwoName{Ass.Prof Supervisor 2} % Navn på vejleder 2
\def\mySupervisorThreeName{Supervisor 3} % Navn på vejleder 3

% Herunder lidt info om din uddannelse og dit institut
\def\myProgramme{Computer Engineering}
\def\myProgrammeDA{Computerteknologi}
\def\myGroupNumber{151}
\def\myDepartment{Department of Electronic Systems}
\def\myDepartmentDA{Institut for Elektroniske Systemer}
\def\myDepartmentWebsite{https://www.es.aau.dk}

% Colophon kommer til at stå på siden efter forsiden og er en reference
% til software og værktøjer brugt til at lave rapporten (og rapportens indhold)
\def\myColophon{%
Typesetting of this report has been done with Overleaf, a LaTeX writing tool.\\
\noindent
During the project simulations has been run in [instert software].\\
Figures have been created in [example: ]Lucidchart and Adobe Illustrator].
}

% Du er nu færdig med at udfylde DocumentSetup
% Kompiler rapporten og se om alt er som det skal være
% Fejl og mangler kan rapporteres som issues på GitHub
% https://github.com/Peasniped/AAU-LateX-Report-template/issues









%
%%%%% Nedenstående skal der ikke rettes i %%%%%
%

%% mySemesterNumber omdannes til Ordinaltal som bogstavtal og nummer (first og 1st)
\def\mySemesterNumberOrdinal{
\ifthenelse{\mySemesterNumber=1}{first}{
\ifthenelse{\mySemesterNumber=2}{second}{
\ifthenelse{\mySemesterNumber=3}{third}{
\ifthenelse{\mySemesterNumber=4}{fourth}{
\ifthenelse{\mySemesterNumber=5}{fifth}{
\ifthenelse{\mySemesterNumber=6}{sixth}{
\ifthenelse{\mySemesterNumber=7}{seventh}{
\ifthenelse{\mySemesterNumber=8}{eighth}{
\ifthenelse{\mySemesterNumber=9}{ninth}{
\ifthenelse{\mySemesterNumber=10}{tenth}{Fejl med def-mySemesterNumberOrdinal
}}}}}}}}}}}

\def\mySemesterNumberOrdinalNumberTh{
\ifthenelse{\mySemesterNumber=1}{1\textsuperscript{st}}{
\ifthenelse{\mySemesterNumber=2}{2\textsuperscript{nd}}{
\ifthenelse{\mySemesterNumber=3}{3\textsuperscript{rd}}{
\ifthenelse{\mySemesterNumber=4}{4\textsuperscript{th}}{
\ifthenelse{\mySemesterNumber=5}{5\textsuperscript{th}}{
\ifthenelse{\mySemesterNumber=6}{6\textsuperscript{th}}{
\ifthenelse{\mySemesterNumber=7}{7\textsuperscript{th}}{
\ifthenelse{\mySemesterNumber=8}{8\textsuperscript{th}}{
\ifthenelse{\mySemesterNumber=9}{9\textsuperscript{th}}{
\ifthenelse{\mySemesterNumber=10}{10\textsuperscript{th}}{Fejl med def-mySemesterNumberOrdinal
}}}}}}}}}}}

%% myAuthors dannes
\ifthenelse{\myAuthorsNUM=1}{
    \def\myAuthors{\myAuthorOneName}
}{ \ifthenelse{\myAuthorsNUM=2}{
    \def\myAuthors{\myAuthorOneName{}, \myAuthorTwoName}
}{ \ifthenelse{\myAuthorsNUM=3}{
    \def\myAuthors{\myAuthorOneName{}, \myAuthorTwoName{}, \myAuthorThreeName}
}{ \ifthenelse{\myAuthorsNUM=4}{
    \def\myAuthors{\myAuthorOneName{}, \myAuthorTwoName{}, \myAuthorThreeName{}, \myAuthorFourName}
}{ \ifthenelse{\myAuthorsNUM=5}{
    \def\myAuthors{\myAuthorOneName{}, \myAuthorTwoName{}, \myAuthorThreeName{}, \myAuthorFourName{}, \myAuthorFiveName}
}{ \ifthenelse{\myAuthorsNUM=6}{
    \def\myAuthors{\myAuthorOneName{}, \myAuthorTwoName{}, \myAuthorThreeName{}, \myAuthorFourName{}, \myAuthorFiveName{}, \myAuthorSixName}
}{ Fejl med def-myAuthors}}}}}}


%% myAuthorsNL dannes
\ifthenelse{\myAuthorsNUM=1}{
    \def\myAuthorsNL{\myAuthorOneName}
}{ \ifthenelse{\myAuthorsNUM=2}{
    \def\myAuthorsNL{\myAuthorOneName{},\\ \myAuthorTwoName}
}{ \ifthenelse{\myAuthorsNUM=3}{
    \def\myAuthorsNL{\myAuthorOneName{},\\ \myAuthorTwoName{},\\ \myAuthorThreeName}
}{ \ifthenelse{\myAuthorsNUM=4}{
    \def\myAuthorsNL{\myAuthorOneName{},\\ \myAuthorTwoName{},\\ \myAuthorThreeName{},\\ \myAuthorFourName}
}{ \ifthenelse{\myAuthorsNUM=5}{
    \def\myAuthorsNL{\myAuthorOneName{},\\ \myAuthorTwoName{},\\ \myAuthorThreeName{},\\ \myAuthorFourName{},\\ \myAuthorFiveName}
}{ \ifthenelse{\myAuthorsNUM=6}{
    \def\myAuthorsNL{\myAuthorOneName{},\\ \myAuthorTwoName{},\\ \myAuthorThreeName{},\\ \myAuthorFourName{},\\ \myAuthorFiveName{},\\ \myAuthorSixName}
}{ Fejl med def-myAuthorsNL}}}}}}


%% mySupervisorsNL dannes
\ifthenelse{\mySupervisorsNUM=1}{
    \def\mySupervisorsNL{\mySupervisorOneName}
}{ \ifthenelse{\mySupervisorsNUM=2}{
    \def\mySupervisorsNL{\mySupervisorOneName{},\\ \mySupervisorTwoName}
}{ \ifthenelse{\mySupervisorsNUM=3}{
    \def\mySupervisorsNL{\mySupervisorOneName{},\\ \mySupervisorTwoName{},\\ \mySupervisorThreeName}
}{ Fejl med def-mySupervisorsNL}}}


%% myThankOurSupervisors dannes
\ifthenelse{\mySupervisorsNUM=1}{
    \def\myThankOurSupervisors{extend our thanks to our supervisor \mySupervisorOneName{} for \mySupervisorOnePronoun{} help and guidance throughout the project.}    
}{ \ifthenelse{\mySupervisorsNUM=2}{
    \def\myThankOurSupervisors{extend our thanks to our supervisors \mySupervisorOneName{} and \mySupervisorTwoName{} for their help and guidance throughout the project.}
}{ \ifthenelse{\mySupervisorsNUM=3}{
    \def\myThankOurSupervisors{extend our thanks to our supervisors \mySupervisorOneName{}, \mySupervisorTwoName{} and \mySupervisorThreeName{} for their help and guidance throughout the project.}
}{ Fejl med def-myThankOurSupervisors}}}% Indlæs opsætning fra DocumentSetup
\pdfbookmark[0]{Front page}{label:frontpage}%
\begin{titlepage}
\newgeometry{top=0cm,bottom=1.2cm,right=0cm,left=0cm}
  
  \thiswatermark{\centering \put(0,-870){\includegraphics[scale=1.0]{media/AAUgraphics/aau_waves.pdf}} }
  % NB: Det ligner i preview, at de der streger er brudte, og generelt i dårlig kvalitet. Det følger ikke med, når man eksporterer PFD'en, der ser de flotte ud.

  \begin{center} %%please do not change the height or width of the frontpage image
    \includegraphics[totalheight=0.5\paperwidth,width=1\paperwidth]{\myForsideBillede}
  \end{center}
	\vspace*{-0.96cm}
  {\noindent\color{aaublue}\fboxsep0pt\colorbox{white}{\begin{tabular}{@{}p{\paperwidth}@{}}
    \centerline{
    \begin{minipage}{0.85\textwidth}
        \bigskip
				\bigskip
        \centering
        \Huge{\textbf{\myProjectTitle}}
    \end{minipage}
    }
	\centerline{
	\begin{minipage}{0.9\textwidth}
        \bigskip
        \centering
        \Large{\mySubtitle}
    \end{minipage}
    }
	\centerline{
	\begin{minipage}{0.8\textwidth}
        \bigskip
        \centering
        {\Large{\textbf{
    \myAuthors{}
        }}}
    \end{minipage}
    }
    \centerline{
    \begin{minipage}{0.9\textwidth}
        \bigskip
        \centering
    \large{\myProgramme{}, \mySemesterNumberOrdinalNumberTh{} semester, Group \myGroupNumber{}\\
    \vspace{0.3cm}
    \mySemester{} \mySemesterYear{}}
    
    \end{minipage}
    }
    \centerline{
    \begin{minipage}{0.9\textwidth}
        \bigskip
        \centering
        \Large{\myProjectNumber}
        \smallskip
    \end{minipage}
    }
  \end{tabular}}}
  \vfill
  \begin{figure}[!b]
	\centering
    \includegraphics[width=0.2\paperwidth]{media/AAUgraphics/aau_logo_circle_en.pdf}
    \vspace{2.5cm}
  \end{figure}
\end{titlepage}
\restoregeometry}{%%
%% Denne fil bruges til at sætte variabler som bruges forskellige steder i rapporten
%% Hvis der er et felt man ikke ønsker at benytte skal det efterlades blankt
%% fx myDenneVariabelBrugesIkke{} - altså uden indhold mellem { }
%%
%% Det kan være nødvendigt at gennemgå de forskellige sider, for at fjerne referencer
%% til udkommenterede variabler, ellers kan resultatet se lidt underligt ud.
%%

%%% Formatering af sideopsætning
\def\myForsideMedBillede{1} % Skal forsiden være med billede? 1 = ja, 0 = nej
\def\myForsideBillede{media/AAUgraphics/frontpageimage.jpg} % Formateres i forholdet(højde:bredde) 1:2
\def\myPageOptimizeForPrint{1} % Skal dokumentet opsættes til print, hvor hvert nyt kapitel er på en lige side i forhold til starten? (Giver mange tomme sider)
% 1 = ja, opsæt til print / 0 = nej, opsæt til web/pdf
\def\myParagraphIndentDisable{1} % Skal indrykning efter afsnit deaktiveres?
% 1 = ja, 0 = nej
\def\myParagraphIndentLength{25pt} % Hvis nej, hvad skal indrykningen være?

% Sidemarginer kan desværre ikke opsættes her, de skal opsættes i "Preamble.tex"
% De findes under punktet 


%%% Tekst på forside/titlepages/colophon/preface
\def\myProjectTitle{Project Title: A Project}
\def\mySubtitle{A project that does sciency things}
\def\myProjectTheme{Theme of project}
\def\mySemester{Fall Semester} % 'Fall semester' eller 'Spring semester'
\def\mySemesterYear{2022}
\def\mySemesterNumber{1} % Hvilket semester er du på? (1,2,3,4,5,6,7,8,9,10)
\def\myProjectNumber{P0 Project} % Hvilket projekt er dette, fx P0 project eller Bachelor's Project
\def\myProjectECTS{15} % Hvor mange ECTS-point er projektet nomeret til?

% Tekst på danske title pages
\def\myProjectTitleDA{Projekttitel: Et projekt}
\def\myProjectThemeDA{Projekttema}
\def\mySemesterDA{Efterårssemester} % 'Forårssemesteret' eller 'Efterårssemesteret'

\def\myAuthorsNUM{4} % Hvor mange gruppemedlemmer der er (1-6)
% Herunder: Navne og mails på gruppemedlemmer
% hvis der fx kun er 4 i gruppen skal nummer 5 og 6 blot være tomme

\def\myAuthorOneName{Author 1} % Navn på gruppemedlem 1
\def\myAuthorTwoName{Author 2} % Navn på gruppemedlem 2
\def\myAuthorThreeName{Author 3} % Navn på gruppemedlem 3
\def\myAuthorFourName{Author 4} % Navn på gruppemedlem 4
\def\myAuthorFiveName{Author 5} % Navn på gruppemedlem 5
\def\myAuthorSixName{Author 6} % Navn på gruppemedlem 6

\def\myAuthorOneEmail{email-1@student.aau.dk} % Email på gruppemedlem 1
\def\myAuthorTwoEmail{email-2@student.aau.dk} % Email på gruppemedlem 2
\def\myAuthorThreeEmail{email-3@student.aau.dk} % Email på gruppemedlem 3
\def\myAuthorFourEmail{email-4@student.aau.dk} % Email på gruppemedlem 4
\def\myAuthorFiveEmail{email-5@student.aau.dk} % Email på gruppemedlem 5
\def\myAuthorSixEmail{email-6@student.aau.dk} % Email på gruppemedlem 6

\def\mySupervisorsNUM{2} % Hvor mange vejledere der er tilknyttet projektet? (1-3)
% Herunder: Navne på Vejledere
% hvis der fx kun er én vejleder tilknyttet projektet efterlades nummer 2 og 3 tomme

\def\mySupervisorOneName{Prof Supervisor 1} % Navn på vejleder 1
\def\mySupervisorOnePronoun{his} % Hvis der kun er én  vejleder, skriv his/her

\def\mySupervisorTwoName{Ass.Prof Supervisor 2} % Navn på vejleder 2
\def\mySupervisorThreeName{Supervisor 3} % Navn på vejleder 3

% Herunder lidt info om din uddannelse og dit institut
\def\myProgramme{Computer Engineering}
\def\myProgrammeDA{Computerteknologi}
\def\myGroupNumber{151}
\def\myDepartment{Department of Electronic Systems}
\def\myDepartmentDA{Institut for Elektroniske Systemer}
\def\myDepartmentWebsite{https://www.es.aau.dk}

% Colophon kommer til at stå på siden efter forsiden og er en reference
% til software og værktøjer brugt til at lave rapporten (og rapportens indhold)
\def\myColophon{%
Typesetting of this report has been done with Overleaf, a LaTeX writing tool.\\
\noindent
During the project simulations has been run in [instert software].\\
Figures have been created in [example: ]Lucidchart and Adobe Illustrator].
}

% Du er nu færdig med at udfylde DocumentSetup
% Kompiler rapporten og se om alt er som det skal være
% Fejl og mangler kan rapporteres som issues på GitHub
% https://github.com/Peasniped/AAU-LateX-Report-template/issues









%
%%%%% Nedenstående skal der ikke rettes i %%%%%
%

%% mySemesterNumber omdannes til Ordinaltal som bogstavtal og nummer (first og 1st)
\def\mySemesterNumberOrdinal{
\ifthenelse{\mySemesterNumber=1}{first}{
\ifthenelse{\mySemesterNumber=2}{second}{
\ifthenelse{\mySemesterNumber=3}{third}{
\ifthenelse{\mySemesterNumber=4}{fourth}{
\ifthenelse{\mySemesterNumber=5}{fifth}{
\ifthenelse{\mySemesterNumber=6}{sixth}{
\ifthenelse{\mySemesterNumber=7}{seventh}{
\ifthenelse{\mySemesterNumber=8}{eighth}{
\ifthenelse{\mySemesterNumber=9}{ninth}{
\ifthenelse{\mySemesterNumber=10}{tenth}{Fejl med def-mySemesterNumberOrdinal
}}}}}}}}}}}

\def\mySemesterNumberOrdinalNumberTh{
\ifthenelse{\mySemesterNumber=1}{1\textsuperscript{st}}{
\ifthenelse{\mySemesterNumber=2}{2\textsuperscript{nd}}{
\ifthenelse{\mySemesterNumber=3}{3\textsuperscript{rd}}{
\ifthenelse{\mySemesterNumber=4}{4\textsuperscript{th}}{
\ifthenelse{\mySemesterNumber=5}{5\textsuperscript{th}}{
\ifthenelse{\mySemesterNumber=6}{6\textsuperscript{th}}{
\ifthenelse{\mySemesterNumber=7}{7\textsuperscript{th}}{
\ifthenelse{\mySemesterNumber=8}{8\textsuperscript{th}}{
\ifthenelse{\mySemesterNumber=9}{9\textsuperscript{th}}{
\ifthenelse{\mySemesterNumber=10}{10\textsuperscript{th}}{Fejl med def-mySemesterNumberOrdinal
}}}}}}}}}}}

%% myAuthors dannes
\ifthenelse{\myAuthorsNUM=1}{
    \def\myAuthors{\myAuthorOneName}
}{ \ifthenelse{\myAuthorsNUM=2}{
    \def\myAuthors{\myAuthorOneName{}, \myAuthorTwoName}
}{ \ifthenelse{\myAuthorsNUM=3}{
    \def\myAuthors{\myAuthorOneName{}, \myAuthorTwoName{}, \myAuthorThreeName}
}{ \ifthenelse{\myAuthorsNUM=4}{
    \def\myAuthors{\myAuthorOneName{}, \myAuthorTwoName{}, \myAuthorThreeName{}, \myAuthorFourName}
}{ \ifthenelse{\myAuthorsNUM=5}{
    \def\myAuthors{\myAuthorOneName{}, \myAuthorTwoName{}, \myAuthorThreeName{}, \myAuthorFourName{}, \myAuthorFiveName}
}{ \ifthenelse{\myAuthorsNUM=6}{
    \def\myAuthors{\myAuthorOneName{}, \myAuthorTwoName{}, \myAuthorThreeName{}, \myAuthorFourName{}, \myAuthorFiveName{}, \myAuthorSixName}
}{ Fejl med def-myAuthors}}}}}}


%% myAuthorsNL dannes
\ifthenelse{\myAuthorsNUM=1}{
    \def\myAuthorsNL{\myAuthorOneName}
}{ \ifthenelse{\myAuthorsNUM=2}{
    \def\myAuthorsNL{\myAuthorOneName{},\\ \myAuthorTwoName}
}{ \ifthenelse{\myAuthorsNUM=3}{
    \def\myAuthorsNL{\myAuthorOneName{},\\ \myAuthorTwoName{},\\ \myAuthorThreeName}
}{ \ifthenelse{\myAuthorsNUM=4}{
    \def\myAuthorsNL{\myAuthorOneName{},\\ \myAuthorTwoName{},\\ \myAuthorThreeName{},\\ \myAuthorFourName}
}{ \ifthenelse{\myAuthorsNUM=5}{
    \def\myAuthorsNL{\myAuthorOneName{},\\ \myAuthorTwoName{},\\ \myAuthorThreeName{},\\ \myAuthorFourName{},\\ \myAuthorFiveName}
}{ \ifthenelse{\myAuthorsNUM=6}{
    \def\myAuthorsNL{\myAuthorOneName{},\\ \myAuthorTwoName{},\\ \myAuthorThreeName{},\\ \myAuthorFourName{},\\ \myAuthorFiveName{},\\ \myAuthorSixName}
}{ Fejl med def-myAuthorsNL}}}}}}


%% mySupervisorsNL dannes
\ifthenelse{\mySupervisorsNUM=1}{
    \def\mySupervisorsNL{\mySupervisorOneName}
}{ \ifthenelse{\mySupervisorsNUM=2}{
    \def\mySupervisorsNL{\mySupervisorOneName{},\\ \mySupervisorTwoName}
}{ \ifthenelse{\mySupervisorsNUM=3}{
    \def\mySupervisorsNL{\mySupervisorOneName{},\\ \mySupervisorTwoName{},\\ \mySupervisorThreeName}
}{ Fejl med def-mySupervisorsNL}}}


%% myThankOurSupervisors dannes
\ifthenelse{\mySupervisorsNUM=1}{
    \def\myThankOurSupervisors{extend our thanks to our supervisor \mySupervisorOneName{} for \mySupervisorOnePronoun{} help and guidance throughout the project.}    
}{ \ifthenelse{\mySupervisorsNUM=2}{
    \def\myThankOurSupervisors{extend our thanks to our supervisors \mySupervisorOneName{} and \mySupervisorTwoName{} for their help and guidance throughout the project.}
}{ \ifthenelse{\mySupervisorsNUM=3}{
    \def\myThankOurSupervisors{extend our thanks to our supervisors \mySupervisorOneName{}, \mySupervisorTwoName{} and \mySupervisorThreeName{} for their help and guidance throughout the project.}
}{ Fejl med def-myThankOurSupervisors}}}% Indlæs opsætning fra DocumentSetup
\pdfbookmark[0]{Front page}{label:frontpage}%
\begin{titlepage}
\vspace*{\fill}
    \backgroundsetup{
    scale=1.1,
    angle=0,
    opacity=1.0,  %% adjust
    contents={\includegraphics[width=\paperwidth,height=\paperheight]{media/AAUgraphics/aau_waves}}
    }
  \addtolength{\hoffset}{0.5\evensidemargin-0.5\oddsidemargin} %set equal margins on the frontpage - remove this line if you want default margins
  \noindent%
  {\color{white}\fboxsep7pt\colorbox{aaublue}{\begin{tabular}{@{}p{\textwidth}@{}}
    \begin{center}
    \Huge{\textbf{\myProjectTitle}}
    \end{center}
    \begin{center}
      \Large{\mySubtitle}
    \end{center}
    \vspace{0.15cm}
   \begin{center}
    \Large{\myAuthors}\\
    \vspace{0.75cm}
    \large{\myProgramme{}, \mySemesterNumber{}. semester, Group \myGroupNumber{}\\
    \vspace{0.3cm}
    \mySemester{} \mySemesterYear{}}
   \end{center}
%% Comment this section in if you are doing Bachelor or Master Project   
   \begin{center}
    \Large{\myProjectNumber}
   \end{center}
  \end{tabular}}}
  \vfill
  \begin{center}
    \includegraphics[width=0.2\paperwidth]{media/AAUgraphics/aau_logo_circle_en}% comment this line in for English version
  \end{center}
\end{titlepage}
\clearpage
}
%%
%% Denne fil bruges til at sætte variabler som bruges forskellige steder i rapporten
%% Hvis der er et felt man ikke ønsker at benytte skal det efterlades blankt
%% fx myDenneVariabelBrugesIkke{} - altså uden indhold mellem { }
%%
%% Det kan være nødvendigt at gennemgå de forskellige sider, for at fjerne referencer
%% til udkommenterede variabler, ellers kan resultatet se lidt underligt ud.
%%

%%% Formatering af sideopsætning
\def\myForsideMedBillede{1} % Skal forsiden være med billede? 1 = ja, 0 = nej
\def\myForsideBillede{media/AAUgraphics/frontpageimage.jpg} % Formateres i forholdet(højde:bredde) 1:2
\def\myPageOptimizeForPrint{1} % Skal dokumentet opsættes til print, hvor hvert nyt kapitel er på en lige side i forhold til starten? (Giver mange tomme sider)
% 1 = ja, opsæt til print / 0 = nej, opsæt til web/pdf
\def\myParagraphIndentDisable{1} % Skal indrykning efter afsnit deaktiveres?
% 1 = ja, 0 = nej
\def\myParagraphIndentLength{25pt} % Hvis nej, hvad skal indrykningen være?

% Sidemarginer kan desværre ikke opsættes her, de skal opsættes i "Preamble.tex"
% De findes under punktet 


%%% Tekst på forside/titlepages/colophon/preface
\def\myProjectTitle{Project Title: A Project}
\def\mySubtitle{A project that does sciency things}
\def\myProjectTheme{Theme of project}
\def\mySemester{Fall Semester} % 'Fall semester' eller 'Spring semester'
\def\mySemesterYear{2022}
\def\mySemesterNumber{1} % Hvilket semester er du på? (1,2,3,4,5,6,7,8,9,10)
\def\myProjectNumber{P0 Project} % Hvilket projekt er dette, fx P0 project eller Bachelor's Project
\def\myProjectECTS{15} % Hvor mange ECTS-point er projektet nomeret til?

% Tekst på danske title pages
\def\myProjectTitleDA{Projekttitel: Et projekt}
\def\myProjectThemeDA{Projekttema}
\def\mySemesterDA{Efterårssemester} % 'Forårssemesteret' eller 'Efterårssemesteret'

\def\myAuthorsNUM{4} % Hvor mange gruppemedlemmer der er (1-6)
% Herunder: Navne og mails på gruppemedlemmer
% hvis der fx kun er 4 i gruppen skal nummer 5 og 6 blot være tomme

\def\myAuthorOneName{Author 1} % Navn på gruppemedlem 1
\def\myAuthorTwoName{Author 2} % Navn på gruppemedlem 2
\def\myAuthorThreeName{Author 3} % Navn på gruppemedlem 3
\def\myAuthorFourName{Author 4} % Navn på gruppemedlem 4
\def\myAuthorFiveName{Author 5} % Navn på gruppemedlem 5
\def\myAuthorSixName{Author 6} % Navn på gruppemedlem 6

\def\myAuthorOneEmail{email-1@student.aau.dk} % Email på gruppemedlem 1
\def\myAuthorTwoEmail{email-2@student.aau.dk} % Email på gruppemedlem 2
\def\myAuthorThreeEmail{email-3@student.aau.dk} % Email på gruppemedlem 3
\def\myAuthorFourEmail{email-4@student.aau.dk} % Email på gruppemedlem 4
\def\myAuthorFiveEmail{email-5@student.aau.dk} % Email på gruppemedlem 5
\def\myAuthorSixEmail{email-6@student.aau.dk} % Email på gruppemedlem 6

\def\mySupervisorsNUM{2} % Hvor mange vejledere der er tilknyttet projektet? (1-3)
% Herunder: Navne på Vejledere
% hvis der fx kun er én vejleder tilknyttet projektet efterlades nummer 2 og 3 tomme

\def\mySupervisorOneName{Prof Supervisor 1} % Navn på vejleder 1
\def\mySupervisorOnePronoun{his} % Hvis der kun er én  vejleder, skriv his/her

\def\mySupervisorTwoName{Ass.Prof Supervisor 2} % Navn på vejleder 2
\def\mySupervisorThreeName{Supervisor 3} % Navn på vejleder 3

% Herunder lidt info om din uddannelse og dit institut
\def\myProgramme{Computer Engineering}
\def\myProgrammeDA{Computerteknologi}
\def\myGroupNumber{151}
\def\myDepartment{Department of Electronic Systems}
\def\myDepartmentDA{Institut for Elektroniske Systemer}
\def\myDepartmentWebsite{https://www.es.aau.dk}

% Colophon kommer til at stå på siden efter forsiden og er en reference
% til software og værktøjer brugt til at lave rapporten (og rapportens indhold)
\def\myColophon{%
Typesetting of this report has been done with Overleaf, a LaTeX writing tool.\\
\noindent
During the project simulations has been run in [instert software].\\
Figures have been created in [example: ]Lucidchart and Adobe Illustrator].
}

% Du er nu færdig med at udfylde DocumentSetup
% Kompiler rapporten og se om alt er som det skal være
% Fejl og mangler kan rapporteres som issues på GitHub
% https://github.com/Peasniped/AAU-LateX-Report-template/issues









%
%%%%% Nedenstående skal der ikke rettes i %%%%%
%

%% mySemesterNumber omdannes til Ordinaltal som bogstavtal og nummer (first og 1st)
\def\mySemesterNumberOrdinal{
\ifthenelse{\mySemesterNumber=1}{first}{
\ifthenelse{\mySemesterNumber=2}{second}{
\ifthenelse{\mySemesterNumber=3}{third}{
\ifthenelse{\mySemesterNumber=4}{fourth}{
\ifthenelse{\mySemesterNumber=5}{fifth}{
\ifthenelse{\mySemesterNumber=6}{sixth}{
\ifthenelse{\mySemesterNumber=7}{seventh}{
\ifthenelse{\mySemesterNumber=8}{eighth}{
\ifthenelse{\mySemesterNumber=9}{ninth}{
\ifthenelse{\mySemesterNumber=10}{tenth}{Fejl med def-mySemesterNumberOrdinal
}}}}}}}}}}}

\def\mySemesterNumberOrdinalNumberTh{
\ifthenelse{\mySemesterNumber=1}{1\textsuperscript{st}}{
\ifthenelse{\mySemesterNumber=2}{2\textsuperscript{nd}}{
\ifthenelse{\mySemesterNumber=3}{3\textsuperscript{rd}}{
\ifthenelse{\mySemesterNumber=4}{4\textsuperscript{th}}{
\ifthenelse{\mySemesterNumber=5}{5\textsuperscript{th}}{
\ifthenelse{\mySemesterNumber=6}{6\textsuperscript{th}}{
\ifthenelse{\mySemesterNumber=7}{7\textsuperscript{th}}{
\ifthenelse{\mySemesterNumber=8}{8\textsuperscript{th}}{
\ifthenelse{\mySemesterNumber=9}{9\textsuperscript{th}}{
\ifthenelse{\mySemesterNumber=10}{10\textsuperscript{th}}{Fejl med def-mySemesterNumberOrdinal
}}}}}}}}}}}

%% myAuthors dannes
\ifthenelse{\myAuthorsNUM=1}{
    \def\myAuthors{\myAuthorOneName}
}{ \ifthenelse{\myAuthorsNUM=2}{
    \def\myAuthors{\myAuthorOneName{}, \myAuthorTwoName}
}{ \ifthenelse{\myAuthorsNUM=3}{
    \def\myAuthors{\myAuthorOneName{}, \myAuthorTwoName{}, \myAuthorThreeName}
}{ \ifthenelse{\myAuthorsNUM=4}{
    \def\myAuthors{\myAuthorOneName{}, \myAuthorTwoName{}, \myAuthorThreeName{}, \myAuthorFourName}
}{ \ifthenelse{\myAuthorsNUM=5}{
    \def\myAuthors{\myAuthorOneName{}, \myAuthorTwoName{}, \myAuthorThreeName{}, \myAuthorFourName{}, \myAuthorFiveName}
}{ \ifthenelse{\myAuthorsNUM=6}{
    \def\myAuthors{\myAuthorOneName{}, \myAuthorTwoName{}, \myAuthorThreeName{}, \myAuthorFourName{}, \myAuthorFiveName{}, \myAuthorSixName}
}{ Fejl med def-myAuthors}}}}}}


%% myAuthorsNL dannes
\ifthenelse{\myAuthorsNUM=1}{
    \def\myAuthorsNL{\myAuthorOneName}
}{ \ifthenelse{\myAuthorsNUM=2}{
    \def\myAuthorsNL{\myAuthorOneName{},\\ \myAuthorTwoName}
}{ \ifthenelse{\myAuthorsNUM=3}{
    \def\myAuthorsNL{\myAuthorOneName{},\\ \myAuthorTwoName{},\\ \myAuthorThreeName}
}{ \ifthenelse{\myAuthorsNUM=4}{
    \def\myAuthorsNL{\myAuthorOneName{},\\ \myAuthorTwoName{},\\ \myAuthorThreeName{},\\ \myAuthorFourName}
}{ \ifthenelse{\myAuthorsNUM=5}{
    \def\myAuthorsNL{\myAuthorOneName{},\\ \myAuthorTwoName{},\\ \myAuthorThreeName{},\\ \myAuthorFourName{},\\ \myAuthorFiveName}
}{ \ifthenelse{\myAuthorsNUM=6}{
    \def\myAuthorsNL{\myAuthorOneName{},\\ \myAuthorTwoName{},\\ \myAuthorThreeName{},\\ \myAuthorFourName{},\\ \myAuthorFiveName{},\\ \myAuthorSixName}
}{ Fejl med def-myAuthorsNL}}}}}}


%% mySupervisorsNL dannes
\ifthenelse{\mySupervisorsNUM=1}{
    \def\mySupervisorsNL{\mySupervisorOneName}
}{ \ifthenelse{\mySupervisorsNUM=2}{
    \def\mySupervisorsNL{\mySupervisorOneName{},\\ \mySupervisorTwoName}
}{ \ifthenelse{\mySupervisorsNUM=3}{
    \def\mySupervisorsNL{\mySupervisorOneName{},\\ \mySupervisorTwoName{},\\ \mySupervisorThreeName}
}{ Fejl med def-mySupervisorsNL}}}


%% myThankOurSupervisors dannes
\ifthenelse{\mySupervisorsNUM=1}{
    \def\myThankOurSupervisors{extend our thanks to our supervisor \mySupervisorOneName{} for \mySupervisorOnePronoun{} help and guidance throughout the project.}    
}{ \ifthenelse{\mySupervisorsNUM=2}{
    \def\myThankOurSupervisors{extend our thanks to our supervisors \mySupervisorOneName{} and \mySupervisorTwoName{} for their help and guidance throughout the project.}
}{ \ifthenelse{\mySupervisorsNUM=3}{
    \def\myThankOurSupervisors{extend our thanks to our supervisors \mySupervisorOneName{}, \mySupervisorTwoName{} and \mySupervisorThreeName{} for their help and guidance throughout the project.}
}{ Fejl med def-myThankOurSupervisors}}}% Indlæs opsætning fra DokumentSetup
\def\myAbstractEN{%
% Engelsk Abstract:
This is where we write out abstract for the project
}

\def\myAbstractDA{%
% Dansk resumé:
Det er her vi skriver vores resumé til projektet1-frontpageImage
}% Indlæs opsætning fra DokumentSetup
\pdfbookmark[0]{English title page}{label:titlepage_en}
\aautitlepage{%
  \englishprojectinfo{\myProjectTitle}{\myProjectTheme}{\mySemester{} \mySemesterYear}{\myGroupNumber}{\myAuthorsNL}{\mySupervisors}{%
    1 % number of printed copies
  }{%
    \today % date of completion
  }%
}{%department and address
  \textbf{\myProgramme}\\
  \myDepartment\\
  Aalborg University\\
  \href{\myDepartmentWebsite}{\myDepartmentWebsite}
}{\myAbstractEN}
\cleardoublepage

{\selectlanguage{danish}
\pdfbookmark[0]{Danish title page}{label:titlepage_da}
\aautitlepage{%
  \danishprojectinfo{\myProjectTitleDA}{\myProjectThemeDA}{\mySemesterDA{} \mySemesterYear}{\myGroupNumber}{\myAuthorsNL}{\mySupervisors}{%
    1 % number of printed copies
  }{%
    \today % date of completion
  }%
}{%department and address
  \textbf{\myProgramme}\\
  \myDepartmentDA\\
  Aalborg Universitet\\
  \href{\myDepartmentWebsite}{\myDepartmentWebsite}
}{\myAbstractDA}}
       % Titelsider på engelsk og på dansk
%%
%% Denne fil bruges til at sætte variabler som bruges forskellige steder i rapporten
%% Hvis der er et felt man ikke ønsker at benytte skal det efterlades blankt
%% fx myDenneVariabelBrugesIkke{} - altså uden indhold mellem { }
%%
%% Det kan være nødvendigt at gennemgå de forskellige sider, for at fjerne referencer
%% til udkommenterede variabler, ellers kan resultatet se lidt underligt ud.
%%

%%% Formatering af sideopsætning
\def\myForsideMedBillede{1} % Skal forsiden være med billede? 1 = ja, 0 = nej
\def\myForsideBillede{media/AAUgraphics/frontpageimage.jpg} % Formateres i forholdet(højde:bredde) 1:2
\def\myPageOptimizeForPrint{1} % Skal dokumentet opsættes til print, hvor hvert nyt kapitel er på en lige side i forhold til starten? (Giver mange tomme sider)
% 1 = ja, opsæt til print / 0 = nej, opsæt til web/pdf
\def\myParagraphIndentDisable{1} % Skal indrykning efter afsnit deaktiveres?
% 1 = ja, 0 = nej
\def\myParagraphIndentLength{25pt} % Hvis nej, hvad skal indrykningen være?

% Sidemarginer kan desværre ikke opsættes her, de skal opsættes i "Preamble.tex"
% De findes under punktet 


%%% Tekst på forside/titlepages/colophon/preface
\def\myProjectTitle{Project Title: A Project}
\def\mySubtitle{A project that does sciency things}
\def\myProjectTheme{Theme of project}
\def\mySemester{Fall Semester} % 'Fall semester' eller 'Spring semester'
\def\mySemesterYear{2022}
\def\mySemesterNumber{1} % Hvilket semester er du på? (1,2,3,4,5,6,7,8,9,10)
\def\myProjectNumber{P0 Project} % Hvilket projekt er dette, fx P0 project eller Bachelor's Project
\def\myProjectECTS{15} % Hvor mange ECTS-point er projektet nomeret til?

% Tekst på danske title pages
\def\myProjectTitleDA{Projekttitel: Et projekt}
\def\myProjectThemeDA{Projekttema}
\def\mySemesterDA{Efterårssemester} % 'Forårssemesteret' eller 'Efterårssemesteret'

\def\myAuthorsNUM{4} % Hvor mange gruppemedlemmer der er (1-6)
% Herunder: Navne og mails på gruppemedlemmer
% hvis der fx kun er 4 i gruppen skal nummer 5 og 6 blot være tomme

\def\myAuthorOneName{Author 1} % Navn på gruppemedlem 1
\def\myAuthorTwoName{Author 2} % Navn på gruppemedlem 2
\def\myAuthorThreeName{Author 3} % Navn på gruppemedlem 3
\def\myAuthorFourName{Author 4} % Navn på gruppemedlem 4
\def\myAuthorFiveName{Author 5} % Navn på gruppemedlem 5
\def\myAuthorSixName{Author 6} % Navn på gruppemedlem 6

\def\myAuthorOneEmail{email-1@student.aau.dk} % Email på gruppemedlem 1
\def\myAuthorTwoEmail{email-2@student.aau.dk} % Email på gruppemedlem 2
\def\myAuthorThreeEmail{email-3@student.aau.dk} % Email på gruppemedlem 3
\def\myAuthorFourEmail{email-4@student.aau.dk} % Email på gruppemedlem 4
\def\myAuthorFiveEmail{email-5@student.aau.dk} % Email på gruppemedlem 5
\def\myAuthorSixEmail{email-6@student.aau.dk} % Email på gruppemedlem 6

\def\mySupervisorsNUM{2} % Hvor mange vejledere der er tilknyttet projektet? (1-3)
% Herunder: Navne på Vejledere
% hvis der fx kun er én vejleder tilknyttet projektet efterlades nummer 2 og 3 tomme

\def\mySupervisorOneName{Prof Supervisor 1} % Navn på vejleder 1
\def\mySupervisorOnePronoun{his} % Hvis der kun er én  vejleder, skriv his/her

\def\mySupervisorTwoName{Ass.Prof Supervisor 2} % Navn på vejleder 2
\def\mySupervisorThreeName{Supervisor 3} % Navn på vejleder 3

% Herunder lidt info om din uddannelse og dit institut
\def\myProgramme{Computer Engineering}
\def\myProgrammeDA{Computerteknologi}
\def\myGroupNumber{151}
\def\myDepartment{Department of Electronic Systems}
\def\myDepartmentDA{Institut for Elektroniske Systemer}
\def\myDepartmentWebsite{https://www.es.aau.dk}

% Colophon kommer til at stå på siden efter forsiden og er en reference
% til software og værktøjer brugt til at lave rapporten (og rapportens indhold)
\def\myColophon{%
Typesetting of this report has been done with Overleaf, a LaTeX writing tool.\\
\noindent
During the project simulations has been run in [instert software].\\
Figures have been created in [example: ]Lucidchart and Adobe Illustrator].
}

% Du er nu færdig med at udfylde DocumentSetup
% Kompiler rapporten og se om alt er som det skal være
% Fejl og mangler kan rapporteres som issues på GitHub
% https://github.com/Peasniped/AAU-LateX-Report-template/issues









%
%%%%% Nedenstående skal der ikke rettes i %%%%%
%

%% mySemesterNumber omdannes til Ordinaltal som bogstavtal og nummer (first og 1st)
\def\mySemesterNumberOrdinal{
\ifthenelse{\mySemesterNumber=1}{first}{
\ifthenelse{\mySemesterNumber=2}{second}{
\ifthenelse{\mySemesterNumber=3}{third}{
\ifthenelse{\mySemesterNumber=4}{fourth}{
\ifthenelse{\mySemesterNumber=5}{fifth}{
\ifthenelse{\mySemesterNumber=6}{sixth}{
\ifthenelse{\mySemesterNumber=7}{seventh}{
\ifthenelse{\mySemesterNumber=8}{eighth}{
\ifthenelse{\mySemesterNumber=9}{ninth}{
\ifthenelse{\mySemesterNumber=10}{tenth}{Fejl med def-mySemesterNumberOrdinal
}}}}}}}}}}}

\def\mySemesterNumberOrdinalNumberTh{
\ifthenelse{\mySemesterNumber=1}{1\textsuperscript{st}}{
\ifthenelse{\mySemesterNumber=2}{2\textsuperscript{nd}}{
\ifthenelse{\mySemesterNumber=3}{3\textsuperscript{rd}}{
\ifthenelse{\mySemesterNumber=4}{4\textsuperscript{th}}{
\ifthenelse{\mySemesterNumber=5}{5\textsuperscript{th}}{
\ifthenelse{\mySemesterNumber=6}{6\textsuperscript{th}}{
\ifthenelse{\mySemesterNumber=7}{7\textsuperscript{th}}{
\ifthenelse{\mySemesterNumber=8}{8\textsuperscript{th}}{
\ifthenelse{\mySemesterNumber=9}{9\textsuperscript{th}}{
\ifthenelse{\mySemesterNumber=10}{10\textsuperscript{th}}{Fejl med def-mySemesterNumberOrdinal
}}}}}}}}}}}

%% myAuthors dannes
\ifthenelse{\myAuthorsNUM=1}{
    \def\myAuthors{\myAuthorOneName}
}{ \ifthenelse{\myAuthorsNUM=2}{
    \def\myAuthors{\myAuthorOneName{}, \myAuthorTwoName}
}{ \ifthenelse{\myAuthorsNUM=3}{
    \def\myAuthors{\myAuthorOneName{}, \myAuthorTwoName{}, \myAuthorThreeName}
}{ \ifthenelse{\myAuthorsNUM=4}{
    \def\myAuthors{\myAuthorOneName{}, \myAuthorTwoName{}, \myAuthorThreeName{}, \myAuthorFourName}
}{ \ifthenelse{\myAuthorsNUM=5}{
    \def\myAuthors{\myAuthorOneName{}, \myAuthorTwoName{}, \myAuthorThreeName{}, \myAuthorFourName{}, \myAuthorFiveName}
}{ \ifthenelse{\myAuthorsNUM=6}{
    \def\myAuthors{\myAuthorOneName{}, \myAuthorTwoName{}, \myAuthorThreeName{}, \myAuthorFourName{}, \myAuthorFiveName{}, \myAuthorSixName}
}{ Fejl med def-myAuthors}}}}}}


%% myAuthorsNL dannes
\ifthenelse{\myAuthorsNUM=1}{
    \def\myAuthorsNL{\myAuthorOneName}
}{ \ifthenelse{\myAuthorsNUM=2}{
    \def\myAuthorsNL{\myAuthorOneName{},\\ \myAuthorTwoName}
}{ \ifthenelse{\myAuthorsNUM=3}{
    \def\myAuthorsNL{\myAuthorOneName{},\\ \myAuthorTwoName{},\\ \myAuthorThreeName}
}{ \ifthenelse{\myAuthorsNUM=4}{
    \def\myAuthorsNL{\myAuthorOneName{},\\ \myAuthorTwoName{},\\ \myAuthorThreeName{},\\ \myAuthorFourName}
}{ \ifthenelse{\myAuthorsNUM=5}{
    \def\myAuthorsNL{\myAuthorOneName{},\\ \myAuthorTwoName{},\\ \myAuthorThreeName{},\\ \myAuthorFourName{},\\ \myAuthorFiveName}
}{ \ifthenelse{\myAuthorsNUM=6}{
    \def\myAuthorsNL{\myAuthorOneName{},\\ \myAuthorTwoName{},\\ \myAuthorThreeName{},\\ \myAuthorFourName{},\\ \myAuthorFiveName{},\\ \myAuthorSixName}
}{ Fejl med def-myAuthorsNL}}}}}}


%% mySupervisorsNL dannes
\ifthenelse{\mySupervisorsNUM=1}{
    \def\mySupervisorsNL{\mySupervisorOneName}
}{ \ifthenelse{\mySupervisorsNUM=2}{
    \def\mySupervisorsNL{\mySupervisorOneName{},\\ \mySupervisorTwoName}
}{ \ifthenelse{\mySupervisorsNUM=3}{
    \def\mySupervisorsNL{\mySupervisorOneName{},\\ \mySupervisorTwoName{},\\ \mySupervisorThreeName}
}{ Fejl med def-mySupervisorsNL}}}


%% myThankOurSupervisors dannes
\ifthenelse{\mySupervisorsNUM=1}{
    \def\myThankOurSupervisors{extend our thanks to our supervisor \mySupervisorOneName{} for \mySupervisorOnePronoun{} help and guidance throughout the project.}    
}{ \ifthenelse{\mySupervisorsNUM=2}{
    \def\myThankOurSupervisors{extend our thanks to our supervisors \mySupervisorOneName{} and \mySupervisorTwoName{} for their help and guidance throughout the project.}
}{ \ifthenelse{\mySupervisorsNUM=3}{
    \def\myThankOurSupervisors{extend our thanks to our supervisors \mySupervisorOneName{}, \mySupervisorTwoName{} and \mySupervisorThreeName{} for their help and guidance throughout the project.}
}{ Fejl med def-myThankOurSupervisors}}}% Indlæs opsætning fra DocumentSetup
\thispagestyle{empty}
{\small
\strut\vfill % push the content to the bottom of the page
\noindent Copyright \copyright{}\mySemesterYear{} Aalborg University\\
\par
\noindent This report is based on a LaTeX template by Morten Zink Stage\\
\url{https://github.com/Peasniped/AAU-LateX-Report-template}.\\
Thanks to Jesper Kjær Nielsen for the work the template was initially adapted from.\par
\vspace{0.2cm}
\noindent \myColophon
}

\clearpage

\pdfbookmark[0]{Contents}{label:contents}
\pagestyle{fancy} %enable headers and footers again
\tableofcontents
\setlength{\parskip}{\myParagraphSkipLength} % Ret i documentSetup - Afstand ml. paragraffer
%%
%% Denne fil bruges til at sætte variabler som bruges forskellige steder i rapporten
%% Hvis der er et felt man ikke ønsker at benytte skal det efterlades blankt
%% fx myDenneVariabelBrugesIkke{} - altså uden indhold mellem { }
%%
%% Det kan være nødvendigt at gennemgå de forskellige sider, for at fjerne referencer
%% til udkommenterede variabler, ellers kan resultatet se lidt underligt ud.
%%

%%% Formatering af sideopsætning
\def\myForsideMedBillede{1} % Skal forsiden være med billede? 1 = ja, 0 = nej
\def\myForsideBillede{media/AAUgraphics/frontpageimage.jpg} % Formateres i forholdet(højde:bredde) 1:2
\def\myPageOptimizeForPrint{1} % Skal dokumentet opsættes til print, hvor hvert nyt kapitel er på en lige side i forhold til starten? (Giver mange tomme sider)
% 1 = ja, opsæt til print / 0 = nej, opsæt til web/pdf
\def\myParagraphIndentDisable{1} % Skal indrykning efter afsnit deaktiveres?
% 1 = ja, 0 = nej
\def\myParagraphIndentLength{25pt} % Hvis nej, hvad skal indrykningen være?

% Sidemarginer kan desværre ikke opsættes her, de skal opsættes i "Preamble.tex"
% De findes under punktet 


%%% Tekst på forside/titlepages/colophon/preface
\def\myProjectTitle{Project Title: A Project}
\def\mySubtitle{A project that does sciency things}
\def\myProjectTheme{Theme of project}
\def\mySemester{Fall Semester} % 'Fall semester' eller 'Spring semester'
\def\mySemesterYear{2022}
\def\mySemesterNumber{1} % Hvilket semester er du på? (1,2,3,4,5,6,7,8,9,10)
\def\myProjectNumber{P0 Project} % Hvilket projekt er dette, fx P0 project eller Bachelor's Project
\def\myProjectECTS{15} % Hvor mange ECTS-point er projektet nomeret til?

% Tekst på danske title pages
\def\myProjectTitleDA{Projekttitel: Et projekt}
\def\myProjectThemeDA{Projekttema}
\def\mySemesterDA{Efterårssemester} % 'Forårssemesteret' eller 'Efterårssemesteret'

\def\myAuthorsNUM{4} % Hvor mange gruppemedlemmer der er (1-6)
% Herunder: Navne og mails på gruppemedlemmer
% hvis der fx kun er 4 i gruppen skal nummer 5 og 6 blot være tomme

\def\myAuthorOneName{Author 1} % Navn på gruppemedlem 1
\def\myAuthorTwoName{Author 2} % Navn på gruppemedlem 2
\def\myAuthorThreeName{Author 3} % Navn på gruppemedlem 3
\def\myAuthorFourName{Author 4} % Navn på gruppemedlem 4
\def\myAuthorFiveName{Author 5} % Navn på gruppemedlem 5
\def\myAuthorSixName{Author 6} % Navn på gruppemedlem 6

\def\myAuthorOneEmail{email-1@student.aau.dk} % Email på gruppemedlem 1
\def\myAuthorTwoEmail{email-2@student.aau.dk} % Email på gruppemedlem 2
\def\myAuthorThreeEmail{email-3@student.aau.dk} % Email på gruppemedlem 3
\def\myAuthorFourEmail{email-4@student.aau.dk} % Email på gruppemedlem 4
\def\myAuthorFiveEmail{email-5@student.aau.dk} % Email på gruppemedlem 5
\def\myAuthorSixEmail{email-6@student.aau.dk} % Email på gruppemedlem 6

\def\mySupervisorsNUM{2} % Hvor mange vejledere der er tilknyttet projektet? (1-3)
% Herunder: Navne på Vejledere
% hvis der fx kun er én vejleder tilknyttet projektet efterlades nummer 2 og 3 tomme

\def\mySupervisorOneName{Prof Supervisor 1} % Navn på vejleder 1
\def\mySupervisorOnePronoun{his} % Hvis der kun er én  vejleder, skriv his/her

\def\mySupervisorTwoName{Ass.Prof Supervisor 2} % Navn på vejleder 2
\def\mySupervisorThreeName{Supervisor 3} % Navn på vejleder 3

% Herunder lidt info om din uddannelse og dit institut
\def\myProgramme{Computer Engineering}
\def\myProgrammeDA{Computerteknologi}
\def\myGroupNumber{151}
\def\myDepartment{Department of Electronic Systems}
\def\myDepartmentDA{Institut for Elektroniske Systemer}
\def\myDepartmentWebsite{https://www.es.aau.dk}

% Colophon kommer til at stå på siden efter forsiden og er en reference
% til software og værktøjer brugt til at lave rapporten (og rapportens indhold)
\def\myColophon{%
Typesetting of this report has been done with Overleaf, a LaTeX writing tool.\\
\noindent
During the project simulations has been run in [instert software].\\
Figures have been created in [example: ]Lucidchart and Adobe Illustrator].
}

% Du er nu færdig med at udfylde DocumentSetup
% Kompiler rapporten og se om alt er som det skal være
% Fejl og mangler kan rapporteres som issues på GitHub
% https://github.com/Peasniped/AAU-LateX-Report-template/issues









%
%%%%% Nedenstående skal der ikke rettes i %%%%%
%

%% mySemesterNumber omdannes til Ordinaltal som bogstavtal og nummer (first og 1st)
\def\mySemesterNumberOrdinal{
\ifthenelse{\mySemesterNumber=1}{first}{
\ifthenelse{\mySemesterNumber=2}{second}{
\ifthenelse{\mySemesterNumber=3}{third}{
\ifthenelse{\mySemesterNumber=4}{fourth}{
\ifthenelse{\mySemesterNumber=5}{fifth}{
\ifthenelse{\mySemesterNumber=6}{sixth}{
\ifthenelse{\mySemesterNumber=7}{seventh}{
\ifthenelse{\mySemesterNumber=8}{eighth}{
\ifthenelse{\mySemesterNumber=9}{ninth}{
\ifthenelse{\mySemesterNumber=10}{tenth}{Fejl med def-mySemesterNumberOrdinal
}}}}}}}}}}}

\def\mySemesterNumberOrdinalNumberTh{
\ifthenelse{\mySemesterNumber=1}{1\textsuperscript{st}}{
\ifthenelse{\mySemesterNumber=2}{2\textsuperscript{nd}}{
\ifthenelse{\mySemesterNumber=3}{3\textsuperscript{rd}}{
\ifthenelse{\mySemesterNumber=4}{4\textsuperscript{th}}{
\ifthenelse{\mySemesterNumber=5}{5\textsuperscript{th}}{
\ifthenelse{\mySemesterNumber=6}{6\textsuperscript{th}}{
\ifthenelse{\mySemesterNumber=7}{7\textsuperscript{th}}{
\ifthenelse{\mySemesterNumber=8}{8\textsuperscript{th}}{
\ifthenelse{\mySemesterNumber=9}{9\textsuperscript{th}}{
\ifthenelse{\mySemesterNumber=10}{10\textsuperscript{th}}{Fejl med def-mySemesterNumberOrdinal
}}}}}}}}}}}

%% myAuthors dannes
\ifthenelse{\myAuthorsNUM=1}{
    \def\myAuthors{\myAuthorOneName}
}{ \ifthenelse{\myAuthorsNUM=2}{
    \def\myAuthors{\myAuthorOneName{}, \myAuthorTwoName}
}{ \ifthenelse{\myAuthorsNUM=3}{
    \def\myAuthors{\myAuthorOneName{}, \myAuthorTwoName{}, \myAuthorThreeName}
}{ \ifthenelse{\myAuthorsNUM=4}{
    \def\myAuthors{\myAuthorOneName{}, \myAuthorTwoName{}, \myAuthorThreeName{}, \myAuthorFourName}
}{ \ifthenelse{\myAuthorsNUM=5}{
    \def\myAuthors{\myAuthorOneName{}, \myAuthorTwoName{}, \myAuthorThreeName{}, \myAuthorFourName{}, \myAuthorFiveName}
}{ \ifthenelse{\myAuthorsNUM=6}{
    \def\myAuthors{\myAuthorOneName{}, \myAuthorTwoName{}, \myAuthorThreeName{}, \myAuthorFourName{}, \myAuthorFiveName{}, \myAuthorSixName}
}{ Fejl med def-myAuthors}}}}}}


%% myAuthorsNL dannes
\ifthenelse{\myAuthorsNUM=1}{
    \def\myAuthorsNL{\myAuthorOneName}
}{ \ifthenelse{\myAuthorsNUM=2}{
    \def\myAuthorsNL{\myAuthorOneName{},\\ \myAuthorTwoName}
}{ \ifthenelse{\myAuthorsNUM=3}{
    \def\myAuthorsNL{\myAuthorOneName{},\\ \myAuthorTwoName{},\\ \myAuthorThreeName}
}{ \ifthenelse{\myAuthorsNUM=4}{
    \def\myAuthorsNL{\myAuthorOneName{},\\ \myAuthorTwoName{},\\ \myAuthorThreeName{},\\ \myAuthorFourName}
}{ \ifthenelse{\myAuthorsNUM=5}{
    \def\myAuthorsNL{\myAuthorOneName{},\\ \myAuthorTwoName{},\\ \myAuthorThreeName{},\\ \myAuthorFourName{},\\ \myAuthorFiveName}
}{ \ifthenelse{\myAuthorsNUM=6}{
    \def\myAuthorsNL{\myAuthorOneName{},\\ \myAuthorTwoName{},\\ \myAuthorThreeName{},\\ \myAuthorFourName{},\\ \myAuthorFiveName{},\\ \myAuthorSixName}
}{ Fejl med def-myAuthorsNL}}}}}}


%% mySupervisorsNL dannes
\ifthenelse{\mySupervisorsNUM=1}{
    \def\mySupervisorsNL{\mySupervisorOneName}
}{ \ifthenelse{\mySupervisorsNUM=2}{
    \def\mySupervisorsNL{\mySupervisorOneName{},\\ \mySupervisorTwoName}
}{ \ifthenelse{\mySupervisorsNUM=3}{
    \def\mySupervisorsNL{\mySupervisorOneName{},\\ \mySupervisorTwoName{},\\ \mySupervisorThreeName}
}{ Fejl med def-mySupervisorsNL}}}


%% myThankOurSupervisors dannes
\ifthenelse{\mySupervisorsNUM=1}{
    \def\myThankOurSupervisors{extend our thanks to our supervisor \mySupervisorOneName{} for \mySupervisorOnePronoun{} help and guidance throughout the project.}    
}{ \ifthenelse{\mySupervisorsNUM=2}{
    \def\myThankOurSupervisors{extend our thanks to our supervisors \mySupervisorOneName{} and \mySupervisorTwoName{} for their help and guidance throughout the project.}
}{ \ifthenelse{\mySupervisorsNUM=3}{
    \def\myThankOurSupervisors{extend our thanks to our supervisors \mySupervisorOneName{}, \mySupervisorTwoName{} and \mySupervisorThreeName{} for their help and guidance throughout the project.}
}{ Fejl med def-myThankOurSupervisors}}}% Indlæs opsætning fra DocumentSetup
\chapter*{Preface\markboth{Preface}{Preface}}\label{ch:preface}
\addcontentsline{toc}{chapter}{Preface}

\noindent
Dette er Forordet(preface) som anvendes til at angive information, der ikke direkte hører hjemme i selve rapporten:\\
 - Hvilken type rapport er dette?\\
 - Hvor er arbejdet udført?\\
 - Tak til personer/virksomheder for assistance i forbindelse med projektet\\
 - Læsevejledning\\
Nogle af disse punkter genereres automatisk ud fra de udfyldte felter i setup/documentSetup. Rapportens forfattere skal desuden underskrive og datere rapporten i bunden af denne side.\\
Kilde: \url{https://www.sdu.dk/-/media/files/om_sdu/institutter/ikbm/diplom+biotek/guide+til+rapportskrivning.pdf}\\
\par
\noindent
This report was written by\mySemesterNumberOrdinal{} semester students (group \myGroupNumber{}), studying \myProgramme{} at Aalborg university during the \mySemester{} of \mySemesterYear{}. The scope of this report is \myProjectECTS{} ECTS credits\\
\par
\noindent
We would like to \myThankOurSupervisors{}\\

\par
\noindent
\textbf{Reading Guide:}
For figures, listings and tables the first number is the chapter and the last numbers are the sequential number of the order they are presented in.\\ 
Source references use a numbered system [Number] with the sources listed in the bibliography section in numerated order at the end of the report.
Acronyms used throughout the report are listed in the glossary in alphabetical order. Each acronym is elucidated where it is first referenced, except for acronyms in tables which are explained in the glossary if necessary.

%% 
%% Vær forsigtig med at rette i nedenstående, da sektionen lidt lever sit eget liv.  
%%

\vspace{\baselineskip}\hfill Aalborg University, \today

\ifthenelse{\myAuthorsNUM=1}{
   \vfill\noindent
    \begin{center}
    \begin{minipage}[b]{0.45\textwidth}
     \centering
     \rule{\textwidth}{0.5pt}
      \myAuthorOneName\\
     {\footnotesize <\myAuthorOneEmail{}>}
    \end{minipage}
    \end{center}

}{ \ifthenelse{\myAuthorsNUM=2}{
    \vfill\noindent
    \begin{minipage}[b]{0.45\textwidth}
     \centering
     \rule{\textwidth}{0.5pt}\\
      \myAuthorOneName\\
     {\footnotesize <\myAuthorOneEmail{}>}
    \end{minipage}
    \hfill
    \begin{minipage}[b]{0.45\textwidth}
     \centering
     \rule{\textwidth}{0.5pt}\\
      \myAuthorTwoName\\
     {\footnotesize <\myAuthorTwoEmail{}>}
    \end{minipage}

}{ \ifthenelse{\myAuthorsNUM=3}{
    \vfill\noindent
    \begin{minipage}[b]{0.45\textwidth}
     \centering
     \rule{\textwidth}{0.5pt}\\
      \myAuthorOneName\\
     {\footnotesize <\myAuthorOneEmail{}>}
    \end{minipage}
    \hfill
    \begin{minipage}[b]{0.45\textwidth}
     \centering
     \rule{\textwidth}{0.5pt}\\
      \myAuthorTwoName\\
     {\footnotesize <\myAuthorTwoEmail{}>}
    \end{minipage}
    
    \vspace{3\baselineskip}
    
    \begin{center}
    \begin{minipage}[b]{0.45\textwidth}
     \centering
     \rule{\textwidth}{0.5pt}
      \myAuthorThreeName\\
     {\footnotesize <\myAuthorThreeEmail{}>}
    \end{minipage}
    \end{center}

}{ \ifthenelse{\myAuthorsNUM=4}{
    \vfill\noindent
    \begin{minipage}[b]{0.45\textwidth}
     \centering
     \rule{\textwidth}{0.5pt}\\
      \myAuthorOneName\\
     {\footnotesize <\myAuthorOneEmail{}>}
    \end{minipage}
    \hfill
    \begin{minipage}[b]{0.45\textwidth}
     \centering
     \rule{\textwidth}{0.5pt}\\
      \myAuthorTwoName\\
     {\footnotesize <\myAuthorTwoEmail{}>}
    \end{minipage}
    
    \vspace{2\baselineskip}
    
    \noindent
    \begin{minipage}[b]{0.45\textwidth}
     \centering
     \rule{\textwidth}{0.5pt}\\
      \myAuthorThreeName\\
     {\footnotesize <\myAuthorThreeEmail{}>}
    \end{minipage}
    \hfill
    \begin{minipage}[b]{0.45\textwidth}
     \centering
     \rule{\textwidth}{0.5pt}\\
      \myAuthorFourName\\
     {\footnotesize <\myAuthorFourEmail{}>}
    \end{minipage}

}{ \ifthenelse{\myAuthorsNUM=5}{
    \vfill\noindent
    \begin{minipage}[b]{0.45\textwidth}
     \centering
     \rule{\textwidth}{0.5pt}\\
      \myAuthorOneName\\
     {\footnotesize <\myAuthorOneEmail{}}
    \end{minipage}
    \vspace{3\baselineskip}
    \hfill
    \begin{minipage}[b]{0.45\textwidth}
     \centering
     \rule{\textwidth}{0.5pt}\\
      \myAuthorTwoName\\
     {\footnotesize <\myAuthorTwoEmail{}>}
    \end{minipage}
    \begin{minipage}[b]{0.45\textwidth}
     \centering
     \rule{\textwidth}{0.5pt}\\
      \myAuthorThreeName\\
     {\footnotesize <\myAuthorThreeEmail{}>}
    \end{minipage}
    \hfill
    \begin{minipage}[b]{0.45\textwidth}
     \centering
     \rule{\textwidth}{0.5pt}\\
      \myAuthorFourName\\
     {\footnotesize <\myAuthorFourEmail{}>}
    \end{minipage}
    \vspace{3\baselineskip}
    \begin{center}
    \begin{minipage}[b]{0.45\textwidth}
     \centering
     \rule{\textwidth}{0.5pt}
      \myAuthorFiveName\\
     {\footnotesize <\myAuthorFiveEmail{}>}
    \end{minipage}
    \end{center}

}{ \ifthenelse{\myAuthorsNUM=6}{
    \vfill\noindent
    \begin{minipage}[b]{0.45\textwidth}
     \centering
     \rule{\textwidth}{0.5pt}\\
      \myAuthorOneName\\
     {\footnotesize <\myAuthorOneEmail{}>}
    \end{minipage}
    \vspace{3\baselineskip}
    \hfill
    \begin{minipage}[b]{0.45\textwidth}
     \centering
     \rule{\textwidth}{0.5pt}\\
      \myAuthorTwoName\\
     {\footnotesize <\myAuthorTwoEmail{}>}
    \end{minipage}
    \begin{minipage}[b]{0.45\textwidth}
     \centering
     \rule{\textwidth}{0.5pt}\\
      \myAuthorThreeName\\
     {\footnotesize <\myAuthorThreeEmail{}>}
    \end{minipage}
    \vspace{3\baselineskip}
    \hfill
    \begin{minipage}[b]{0.45\textwidth}
     \centering
     \rule{\textwidth}{0.5pt}\\
      \myAuthorFourName\\
     {\footnotesize <\myAuthorFourEmail{}>}
    \end{minipage}
    \begin{minipage}[b]{0.45\textwidth}
     \centering
     \rule{\textwidth}{0.5pt}\\
      \myAuthorFiveName\\
     {\footnotesize <\myAuthorFiveEmail{}>}
    \end{minipage}
    \hfill
    \begin{minipage}[b]{0.45\textwidth}
     \centering
     \rule{\textwidth}{0.5pt}\\
      \myAuthorSixName\\
     {\footnotesize <\myAuthorSixEmail{}>}
    \end{minipage}
}{ !!! fejl: myAuthorsNUM skal være et tal mellem 1-6, myAuthorsNUM er \myAuthorsNUM !!! }}}}}}

% Printer Forklaring til forkortelser og ordbog
% acro v2 syntax:
\printacronyms[]
\addcontentsline{toc}{chapter}{Acronyms} 
\printglossaries
\pagebreak

\pagenumbering{arabic} %use Arabic page numbering in the main matter
\vspace{-2cm}
\chapter{Introduction}\label{ch:introduction}
\ifthenelse{\myExpandAcronymsEveryChapter=1}{\acresetall}{} % Nulstiller alle akronymers "er brugt"-tilstand. Indstilles i documentSetup.

Here is the introduction. The next chapter is chapter. ~\ref{ch:ch2label}.

a new paragraph

\section{Examples}
You can also have examples in your document such as in example~\ref{ex:simple_example}.
\begin{example}{An Example of an Example}
    \label{ex:simple_example}
    Here is an example with some math
    \begin{equation}
        0 = \exp(i\pi)+1\ .
    \end{equation}
    You can adjust the colour and the line width in the {\tt macros.tex} file.
    You can also use acronyms with \ac{eksempel}. The second time you refer to it, it is written like this \ac{eksempel}
  
    you can also define terms and explain them in the glossary with \gls{eksempel}.
    \vspace{0.5cm}
\end{example}


\myTitle{Example of figure}
If you want an image in the text this is how to do it:
\begin{figure}[H] % HUSK [H], så bliver billedet der hvor den er sat
    \centering
    \includegraphics[width=0.25\linewidth]{media/Example monki.jpg}
    \caption{Visualization of a \myExpression{Monkey}}
    \label{fig:monki}
\end{figure}


\begin{table}[h!]
    \centering
    \begin{tblr}{
        colspec = {l l l},
        vlines = {},    % Vertikale linjer
    	hlines = {},    % Horisontale linjer
    	rows = {bg = aaublue!15},   % farver alle rækker blå med 15% opacity
    	row{odd} = {bg = aaublue!30},   % farver hver anden række blå med 30% opacity (mørkere)
    	column{1} = {bg = aaublue!65, fg = white,halign = c},    % farver første kolonne blå/lilla og gør teksten hvid og centrerer tekst
    	row{1} = {bg = aaublue, fg = white},    % farver øverste række blå med hvid tekst
    	% merging cells
    	cell{2}{1} = {r = 2}{valign = m},   % merger celle 2,1 en række ned
    	cell{4}{2} = {c = 2}{halign = c},   % merger celle 4,2 en kolonne hen
        }
        \textbf{one} & \textbf{two} & \textbf{three}\\
        four         & five         & six\\
                     & seven        & eight\\
        nine         & ten          & \\
    \end{tblr}
    \caption{this is an example of a table(tabularray / tblr) with colors and merged cells}
    \label{tab:table}
\end{table}

\section{How Does Sections, Subsections, and Subsections Look?}
Well, like this
\subsection{This is a Subsection}
and this
\subsubsection{This is a Subsubsection}
and this.

\paragraph{A Paragraph}
You can also use paragraph titles which look like this.

This is a reference to \autoref{}{Table}{tab:table}
\newline
This is a reference with a page number to \myAutoRefOnPage{Table}{tab:table}

\subparagraph{A Subparagraph} Moreover, you can also use subparagraph titles which look like this\todo{Is it possible to add a subsubparagraph?}. They have a small indentation as opposed to the paragraph titles.

\todo[inline,color=green]{I think that a summary of this exciting chapter should be added.}

\section{The best section}
\label{this:good}

This is a nameref to \nameref{this:good}\\
This is a ref to \ref{this:good}\\
This is an autoref to \autoref{this:good}\\
\setlength{\headheight}{14pt}
\chapter{Chapter 2 name}\label{ch:ch2label}
\ifthenelse{\myExpandAcronymsEveryChapter=1}{\acresetall}{} % Nulstiller alle akronymers "er brugt"-tilstand. Indstilles i documentSetup.

Here is chapter 2. If you want to leearn \todo{I think this word is mispelled} more about \LaTeXe{}, have a look at \cite{Madsen2010}, \cite{Oetiker2010} and \cite{Mittelbach2005}.
\missingfigure{We need a figure right here!}
\chapter{Requirements}\label{ch:requirements}
\ifthenelse{\myExpandAcronymsEveryChapter=1}{\acresetall}{} % Nulstiller alle akronymers "er brugt"-tilstand. Indstilles i documentSetup.
\chapter{Design}\label{ch:design}
\chapter{Implementation}\label{ch:implementation}
\chapter{Test and Validation}\label{ch:test and validation}
\chapter{Conclusion}\label{ch:conclusion}
\ifthenelse{\myExpandAcronymsEveryChapter=1}{\acresetall}{} % Nulstiller alle akronymers "er brugt"-tilstand. Indstilles i documentSetup.

There are probably still some bugs in the template. If you should find one, then please submit it on \url{https://github.com/Peasniped/AAU-LateX-Report-template/issues}.

\printbibliography[heading=bibintoc, title=Bibliography]
\label{bib:bibliography}

\label{myLastPage}
\clearpage

\appendix
%NB: Kommandoen '\myAppendixPageNumbering' kaldes før indsættelsen af et appendix(bilag). Kommandoen nulstiller sidetallet og parameteren 'A-' definerer prefix på sidetallene, så side 11 i bilag A formateres som side 'A-11'. Kommandoen skal kaldes på ny før hvert bilag(fx, bilag A, B, C etc.).
\myAppendixPageNumbering{A-}
\chapter{Appendix A name}\label{ch:appAlabel}
Here is the first appendix

\clearpage

\myAppendixPageNumbering{B-}
\chapter{Kodeeksempler}\label{ch:appendix-kodeeksempler}

\section{Enumerate / Itemize}
Enumerate bruges til at lave en liste, fx med tal, bogstaver eller romertal. Itemize bruges til bullets.

\begin{lstlisting}[language=Tex, caption=Her vises enumerate med forskellige formatering af tallet]
 \begin{enumerate}[label=\textbf{\arabic*: }]
    \item This is the first item
    \item This is the second item
\end{enumerate}

 \begin{enumerate}[label=\textit{\roman*.  }]   % roman = i ii iii, Roman = I II III
    \item This is the third item
    \item This is the fourth item
\end{enumerate}

 \begin{enumerate}[label=\Alph*) ]    % alph = a b c, Alph = A B C
    \item This is the fifth item
    \item This is the sixth item
\end{enumerate}

\begin{itemize}[itemsep=1pt]    % itemsep = afstand ml punkter
    \item This is the final item
\end{itemize}
\end{lstlisting}

\subsection{Eksempel:}

 \begin{enumerate}[label=\textbf{\arabic*: }]
    \item This is the first item
    \item This is the second item
\end{enumerate}

 \begin{enumerate}[label=\textit{\roman*.  }]   % roman = i ii iii, Roman = I II III
    \item This is the third item
    \item This is the fourth item
\end{enumerate}

 \begin{enumerate}[label=\Alph*) ]    % alph = a b c, Alph = A B C
    \item This is the fifth item
    \item This is the sixth item
\end{enumerate}

\begin{itemize}[itemsep=1pt]    % itemsep = afstand ml punkter
    \item This is the final item
\end{itemize}

\pagebreak

\section{wrapTable / wrapFigure}
Wraptable/figure bruges til at få teksten til at gå rundt om billedet. Ligesom i Word, hvor man kan vælge at wrape teksten rundt om et billede.
De to environments har samme formatering, men wrapTable gør så elementet referes som tabel, hvis man bruger en label, og wrapFigure refererer til . 

\begin{lstlisting}[language=Tex, caption=Her vises brug af wrapTable]
    \lipsum[4]    % Dummytekst
    \begin{wrapfigure}{r}{8cm} % Bredden af wrap-containeren
    
        \centering  % Centrerer billedet i wrap-containeren
        \includegraphics[width=0.95\linewidth]{media/AAUgraphics/frontpageImage.jpg} % width er billedets bredde, linewidth refererer til bredden af wrap-containeren
        
        \caption{SPAAAAAAAAAAAAAAAAACEEEEEEEE!}
        \label{fig:wraptable}
        
    \end{wrapfigure}
    
    \lipsum[4-5]    % Dummytekst
\end{lstlisting}

\subsection{Eksempel:}
\lipsum[4]    % Dummytekst
\begin{wrapfigure}{r}{8cm} % Bredden af containeren
    \centering
    \includegraphics[width=0.95\linewidth]{media/AAUgraphics/frontpageImage.jpg} 
    % width er billedets bredde, linewidth refererer til bredden af containeren, så 0.95*linewidth giver lidt margin om billedet. 
    \caption{SPAAAAAAAAAAAAAAAAACEEEEEEEE!}
    \label{fig:wraptable}
\end{wrapfigure}
\lipsum[4-5]    % Dummytekst

\pagebreak

\section{Minipages - To billeder ved siden af hinanden}
Hvis vi gerne vil have to billeder ved siden af hinanden kan vi dele siden op i to minipages.

\begin{lstlisting}[language=Tex, caption=Her vises brug af wrapTable]
\lipsum[64]
\par
\vspace{1cm}
\noindent % vigtigt at have med
\begin{minipage}{0.5\textwidth} % hver minipage er 0.5 sidebredde i bred
    \centering
    \includegraphics[width=0.8\linewidth]{media/AAUgraphics/aau_logo_circle_en.pdf} 
    \captionof{figure}{AAUs cirkel-logo paa engelsk} % NB: captionof ikke caption
    \label{fig:aauLogoEN}
\end{minipage}%
%
\begin{minipage}{0.5\textwidth}
    \centering
    \includegraphics[width=0.8\linewidth]{media/AAUgraphics/aau_logo_circle_da.pdf} 
    \captionof{figure}{AAUs cirkel-logo paa engelsk} % NB: captionof ikke caption
    \label{fig:aauLogoDA}
\end{minipage}
\end{lstlisting}

\noindent
Note: Vi bruger kommandoen captionof{figure}{caption skrives her} i stedet for almindeligvist at bruge caption{caption skrives her}, da vi ikke er inde i et korrekt "float environment". Det er fordi vi er inde i et minipage-environment som ikke definerer om det er en figur eller en tabel der er tale om.

\subsection{Eksempel:}
\lipsum[64]
\par
\vspace{1cm}
\noindent % vigtigt at have med
\begin{minipage}{0.5\textwidth} % hver minipage er 0.5 sidebredde i bred
    \centering
    \includegraphics[width=0.65\linewidth]{media/AAUgraphics/aau_logo_circle_en.pdf} 
    \captionof{figure}{AAUs cirkel-logo på engelsk} % bemærk \captionof ikke \caption
    \label{fig:aauLogoEN}
\end{minipage}%
%
\begin{minipage}{0.5\textwidth}
    \centering
    \includegraphics[width=0.65\linewidth]{media/AAUgraphics/aau_logo_circle_da.pdf} 
    \captionof{figure}{AAUs cirkel-logo på engelsk} % bemærk \captionof ikke \caption
    \label{fig:aauLogoDA}
\end{minipage}





\end{document}